% =====================================================================
%  Tesoreria Distribuida  --  Manual tecnico de ESTUDIO  (Equipo 18)
%  VERSION DE ESTUDIO -- NO ENTREGAR.
%  Documento interno para preparar la explicacion del proyecto ante el
%  profesor. No forma parte de la entrega oficial.
%
%  Compilar:  cd ~/hatziry/docs && ./compilar.sh
%             (o: pdflatex manual-estudio.tex  -- dos veces, por el indice)
% =====================================================================
\documentclass[11pt,a4paper]{article}

\usepackage[utf8]{inputenc}
\usepackage[T1]{fontenc}
\usepackage{amsmath}
\usepackage[spanish,es-noshorthands]{babel}
\usepackage[margin=2.5cm]{geometry}
\usepackage{xcolor}
\usepackage{listings}
\usepackage{graphicx}
\usepackage{enumitem}
\usepackage{fancyhdr}
\usepackage[hidelinks]{hyperref}

% ---- Marcador de contenido pendiente (sale en rojo y compila igual) ----
\newcommand{\pendiente}[1]{{\color{red}\textbf{[PENDIENTE: #1]}}}

% ---- Codigo en linea (clases, metodos, rutas, variables) ----
\newcommand{\codigo}[1]{\texttt{#1}}

% ---- Colores y estilo de listados de codigo ----
\definecolor{codegray}{gray}{0.95}
\definecolor{kw}{RGB}{0,0,180}
\definecolor{cmt}{RGB}{0,128,0}
\definecolor{str}{RGB}{163,21,21}

\lstdefinestyle{tes}{
  backgroundcolor=\color{codegray},
  basicstyle=\ttfamily\footnotesize,
  keywordstyle=\color{kw}\bfseries,
  commentstyle=\color{cmt}\itshape,
  stringstyle=\color{str},
  numbers=left, numberstyle=\tiny\color{gray}, numbersep=8pt,
  showstringspaces=false, breaklines=true, frame=single, tabsize=2,
  captionpos=b,
}
% Lenguaje vacio para listados de JSON/bash/HTTP (la guia usa language=none).
\lstdefinelanguage{none}{keywords={},sensitive=false,morecomment=[l]{},morestring=[b]"}
\lstset{style=tes}

% ---- Encabezado / pie ----
\pagestyle{fancy}\fancyhf{}
\fancyhead[L]{Tesoreria Distribuida -- Manual de estudio}
\fancyhead[R]{\thepage}
\renewcommand{\headrulewidth}{0.4pt}
\fancyfoot[C]{\small\color{red}VERSION DE ESTUDIO -- NO ENTREGAR}

\begin{document}

% ---------------------------- PORTADA -------------------------------
\begin{titlepage}
\centering
\vspace*{3cm}
{\Huge\bfseries Tesoreria Distribuida\par}
\vspace{0.6cm}
{\Large Manual tecnico de estudio\par}
\vspace{0.3cm}
{\large Equipo 18 -- Proyecto Final de Sistemas Distribuidos\par}
\vspace{2.5cm}
{\large\bfseries\color{red}VERSION DE ESTUDIO -- NO ENTREGAR\par}
\vspace{0.4cm}
\begin{minipage}{0.8\textwidth}\centering\normalsize
Documento interno para entender y poder explicar el sistema ante el
profesor. \textbf{No} forma parte de la entrega oficial y puede parecerse
a otra documentacion: es intencional, su unico fin es estudiar.
\end{minipage}
\vfill
{\normalsize Mini banco distribuido en Java puro\\ Paquete \codigo{mx.ipn.escom.tesoreria}\par}
\end{titlepage}

% ---------------------------- INDICE --------------------------------
\tableofcontents
\newpage

% --------------------------- CAPITULOS ------------------------------
% Cada archivo de capitulos/ es un FRAGMENTO que empieza con \section.
% El otro agente los rellena leyendo el codigo real (ver PROMPT-AGENTE.md
% y GUIA.md). Reemplazar cada \pendiente{...} con prosa real.

\section{Arquitectura general}

\subsection{El sistema en una pagina}

\textbf{Tesoreria Distribuida} (Equipo 18) es un mini banco distribuido escrito
en Java puro, sin frameworks web ni servidor de aplicaciones. El estado que
mantiene es un \emph{ledger} de cuentas (un libro de saldos) replicado entre
tres nodos. Sobre ese ledger ofrece cuatro capacidades de cliente: registro de
usuarios, \emph{login} que emite un token JWT, consulta del saldo de una cuenta y
transferencia de dinero entre dos cuentas. A esas operaciones se suman la
\textbf{durabilidad} de la bitacora de transferencias en Cloud Storage (GCS) y un
\textbf{dashboard} de monitoreo del cluster.

El paquete raiz es \codigo{mx.ipn.escom.tesoreria}. El build con Maven produce un
\emph{fat jar} ejecutable (un \codigo{jar-with-dependencies} cuya clase principal
es \codigo{app/Node.java}). Una propiedad central del diseno es que \emph{un mismo
jar} arranca como lider o como replica segun la configuracion: el rol no se
compila, se decide en tiempo de ejecucion leyendo variables de entorno con
prefijo \codigo{TES\_}. En concreto, \codigo{NodeConfig.isLeader()} considera al
proceso lider cuando \codigo{TES\_LEADER\_HOST} esta vacio o ausente, y replica
cuando esa variable nombra el host del lider a seguir. El secreto de firma de los
tokens (\codigo{TES\_JWT\_SECRET}) debe tener el MISMO valor en los tres nodos
para que un JWT emitido en cualquiera sea valido en todos.

Al arrancar, \codigo{Node.main} carga el dataset compartido de cuentas en un
\codigo{core/Ledger.java} nuevo y lo envuelve en un \codigo{core/Bank.java}, que
es quien posee la secuencia de \emph{commits} y la difusion a observadores. Si la
carga no produce cuentas, el nodo falla de inmediato en lugar de servir con un
ledger vacio:

\begin{lstlisting}[language=Java,caption={Node.main: carga del dataset y rechazo de un ledger vacio}]
Ledger ledger = new Ledger();
int loaded = Dataset.loadInto(ledger, Path.of(config.datasetPath()), config.idBase());
if (loaded <= 0) {
    throw new IllegalStateException("dataset loaded 0 accounts from "
            + config.datasetPath() + "; refusing to start with an empty ledger");
}
Bank bank = new Bank(ledger);
\end{lstlisting}

\subsection{Topologia: lider y replicas}

El sistema corre en tres nodos. \textbf{nodo-1} es el \textbf{lider}
(maquina \codigo{e2-standard-4}) y es el unico que recibe el trafico de clientes:
sobre el se ejecutan registro, \emph{login}, consulta de saldo y, sobre todo, las
transferencias, porque es el lider quien aplica de verdad cada movimiento, le
asigna un numero de secuencia y lo numera de forma monotona. \textbf{nodo-2} y
\textbf{nodo-3} son \textbf{replicas} (maquinas \codigo{e2-standard-2}) que no
originan transacciones de cliente: solo siguen el flujo de transferencias que el
lider ya acepto y lo reaplican localmente para converger al mismo estado.

Quien es lider y quien es replica lo decide la configuracion, no el codigo.
\codigo{app/NodeConfig.java} lee \codigo{TES\_LEADER\_HOST} desde el entorno y
expone el rol con \codigo{isLeader()}:

\begin{lstlisting}[language=Java,caption={NodeConfig.isLeader: el rol se infiere de TES\_LEADER\_HOST}]
public boolean isLeader() {
    return leaderHost == null || leaderHost.isBlank();
}
\end{lstlisting}

A partir de ese rol, \codigo{Node.main} ramifica el arranque. En el \textbf{lider}
(cuando \codigo{config.isLeader()} es verdadero) corre, si GCS esta configurado,
la recuperacion en frio de la bitacora con \codigo{durable/Journal.java}, y
registra dos observadores sobre el banco: el propio \codigo{Journal}, que guarda
de forma durable cada transferencia nueva, y un \codigo{cluster/ReplicaFeed.java},
un servidor TCP que transmite los \emph{commits} en vivo hacia las replicas
conectadas. En la \textbf{replica} (cualquier otro caso) se restaura primero el
\emph{checkpoint} durable con \codigo{durable/ReplicaCheckpoint.java} si hay GCS, y
luego se arranca un \codigo{cluster/ReplicaSync.java} que se pone al dia desde el
lider por numero de secuencia y aplica los \emph{commits} en vivo. El cliente pega
al lider porque solo alli existe el camino que origina, valida y numera una
transferencia; las replicas son un seguidor convergente, no un punto de entrada de
escritura.

\subsection{Mapa de paquetes}

El codigo se organiza por responsabilidad bajo \codigo{mx.ipn.escom.tesoreria}:

\begin{itemize}
  \item \codigo{app} --- arranque y configuracion del proceso: \codigo{Node}
        (el \codigo{main}), \codigo{NodeConfig} (variables \codigo{TES\_*}) y
        \codigo{NodeStats} (metricas del nodo).
  \item \codigo{net} --- transporte HTTP sobre NIO crudo, sin librerias:
        \codigo{Server}, \codigo{IoLoop}, \codigo{Channel},
        \codigo{RequestParser}, \codigo{Request}, \codigo{Reply},
        \codigo{Routes} y la interfaz \codigo{Endpoint}.
  \item \codigo{core} --- el dominio del banco: \codigo{Account},
        \codigo{Ledger} (cuentas y dinero), \codigo{Bank} (orquesta la
        transaccion y la secuencia), \codigo{Money}, \codigo{Transfer},
        \codigo{TransferLog}, \codigo{TransferException}, \codigo{Dataset} y la
        interfaz \codigo{CommitListener}.
  \item \codigo{security} --- autenticacion: \codigo{Tokens} (JWT),
        \codigo{Passwords} (bcrypt), \codigo{Credential},
        \codigo{CredentialStore} y \codigo{Authenticator}.
  \item \codigo{api} --- los \emph{endpoints} HTTP: \codigo{RegisterEndpoint},
        \codigo{LoginEndpoint}, \codigo{BalanceEndpoint},
        \codigo{TransferEndpoint}, \codigo{StatsEndpoint},
        \codigo{PanelEndpoint} y \codigo{DashboardEndpoint}.
  \item \codigo{cluster} --- replicacion TCP (una linea JSON por transferencia):
        \codigo{ReplicaFeed} en el lider, \codigo{ReplicaSync} en la replica y
        \codigo{WireCodec} para codificar el cable.
  \item \codigo{durable} --- durabilidad en GCS por REST: \codigo{Journal},
        \codigo{GcsStore}, \codigo{GcsAuth} y \codigo{ReplicaCheckpoint}.
  \item \codigo{loadtest} --- \codigo{LoadDriver}, generador de carga y
        verificador de la invariante de conservacion del dinero.
\end{itemize}

\subsection{Flujo de una transferencia de punta a punta}

Una transferencia recorre todas las capas en orden. Tomando como ejemplo
\codigo{POST /api/transactions/transfer} sobre el lider:

\begin{enumerate}
  \item \textbf{Cliente $\rightarrow$ red.} La peticion llega al servidor NIO
        \codigo{net/Server.java}. El hilo aceptador la reparte por
        \emph{round-robin} a un reactor (\codigo{net/IoLoop.java}); los reactores
        hacen la E/S y un pool de \emph{workers} computa la respuesta.
  \item \textbf{Ruteo.} \codigo{net/Routes.java} empareja metodo y ruta y
        despacha al \emph{endpoint} registrado en \codigo{Node.main} para
        \codigo{POST /api/transactions/transfer}: una instancia de
        \codigo{api/TransferEndpoint.java}.
  \item \textbf{Validacion JWT.} El \emph{endpoint} primero exige
        \codigo{POST} (devuelve 405 en otro caso) y luego autoriza el token con
        \codigo{security/Authenticator.java}; sin token valido responde 401.
  \item \textbf{Parseo y conversion a centavos.} Lee del cuerpo JSON las claves
        obligatorias \codigo{sourceAccountId}, \codigo{targetAccountId} (ids como
        cadenas) y \codigo{amount} (decimal), convirtiendo el monto a centavos
        con \codigo{core/Money.java}. Un cuerpo malformado da 400
        \codigo{bad\_request}; un id numerico fuera del rango \codigo{int} da 404
        \codigo{no\_such\_account}.
  \item \textbf{Banco y ledger.} Entrega el movimiento a
        \codigo{Bank.transfer(from, to, cents)}, que rechaza por politica la
        auto-transferencia y el monto no positivo, delega el debito/credito atomico
        al \codigo{core/Ledger.java}, y solo si tiene exito estampa el numero de
        secuencia, lo registra en el \codigo{core/TransferLog.java} y avanza los
        contadores.
  \item \textbf{Difusion a observadores.} En el mismo \codigo{transfer}, el banco
        notifica a cada \codigo{CommitListener} registrado: el
        \codigo{durable/Journal.java}, que persiste la transferencia en la
        bitacora GCS, y el \codigo{cluster/ReplicaFeed.java}, que la empuja a las
        replicas conectadas.
  \item \textbf{Feed a replicas.} Cada replica recibe el \emph{commit} por su
        \codigo{cluster/ReplicaSync.java} y lo reaplica con
        \codigo{Bank.applyReplicated}, que es idempotente por secuencia para no
        duplicar saldos tras una reconexion.
  \item \textbf{Respuesta.} El \emph{endpoint} devuelve 200 con
        \codigo{\{"status":"ok","seq":<seq>\}}; un rechazo del banco
        (\codigo{TransferException}) se mapea a 404 para \codigo{no\_such\_account}
        y a 400 para los demas codigos, con cuerpo \codigo{\{"error":<code>\}}.
\end{enumerate}

El nucleo del paso 5, ya en \codigo{api/TransferEndpoint.java}, entrega el
movimiento al banco y traduce el resultado a una respuesta REST:

\begin{lstlisting}[language=Java,caption={TransferEndpoint.serve: del banco a la respuesta HTTP}]
try {
    long seq = bank.transfer(from, to, cents);
    JsonObject ok = new JsonObject();
    ok.addProperty("status", "ok");
    ok.addProperty("seq", Long.valueOf(seq));
    return Reply.json(200, gson.toJson(ok));
} catch (TransferException e) {
    int status = "no_such_account".equals(e.code()) ? 404 : 400;
    JsonObject err = new JsonObject();
    err.addProperty("error", e.code());
    return Reply.json(status, gson.toJson(err));
}
\end{lstlisting}

\section{Arranque y configuracion del nodo}

\subsection{Build: fat jar con Maven}

El proyecto se compila con Maven y produce un \emph{fat jar} ejecutable
descrito en \codigo{pom.xml}. El artefacto es \codigo{tesoreria-distribuida}
version \codigo{1.0.0}, empaquetado como \codigo{jar} y compilado contra Java~17
(\codigo{maven.compiler.release} = 17). El comando habitual es
\codigo{mvn clean package}, que dispara el plugin
\codigo{maven-assembly-plugin} (version 3.7.1) con el descriptor
\codigo{jar-with-dependencies}: ese descriptor reune el codigo del nodo y todas
sus dependencias en un solo jar autocontenido cuyo manifiesto fija la clase
principal \codigo{mx.ipn.escom.tesoreria.app.Node}. Asi un unico jar arranca
indistintamente como lider o como replica, segun las variables de entorno.

\begin{lstlisting}[language=none,caption={Plugin assembly: fat jar y clase principal (pom.xml)}]
<artifactId>maven-assembly-plugin</artifactId>
<version>3.7.1</version>
<configuration>
    <descriptorRefs>
        <descriptorRef>jar-with-dependencies</descriptorRef>
    </descriptorRefs>
    <archive>
        <manifest>
            <mainClass>mx.ipn.escom.tesoreria.app.Node</mainClass>
        </manifest>
    </archive>
</configuration>
\end{lstlisting}

El sistema es Java puro, sin frameworks: la lista de dependencias declaradas se
limita a tres bibliotecas, \codigo{org.mindrot:jbcrypt:0.4} (hash de
contrasenas), \codigo{com.auth0:java-jwt:4.3.0} (emitir y validar JWT) y
\codigo{com.google.code.gson:gson:2.13.1} (parseo de JSON). Es importante para
la rubrica que el \codigo{pom.xml} \emph{no} traiga ningun driver JDBC ni base
de datos embebida: el requisito de usar un SGBD se cumple con cero, porque el
estado vive en memoria (el ledger) y la durabilidad se resuelve contra Cloud
Storage. De hecho el comentario del propio \codigo{pom.xml} aclara que no se
permite el SDK de Google Cloud, y que GCS se alcanza por su API REST JSON cruda
usando \codigo{java.net.http.HttpClient}.

\subsection{Punto de entrada: Node}

La clase \codigo{app/Node.java} es el punto de entrada del proceso. Su
\codigo{main()} recibe opcionalmente el puerto HTTP como \codigo{args[0]} (por
defecto 8080) y orquesta el arranque en una secuencia fija. Primero lee la
configuracion con \codigo{NodeConfig.fromEnv()}; despues carga el dataset de
cuentas en un \codigo{Ledger} nuevo con
\codigo{Dataset.loadInto(ledger, ...)} y, si la carga devuelve cero cuentas,
falla de inmediato con \codigo{IllegalStateException} en vez de servir un ledger
vacio (sintoma tipico de un CSV ausente o truncado). El ledger ya poblado se
envuelve en un \codigo{Bank}, que es quien posee la secuencia de \emph{commits}
y el reparto a los \codigo{CommitListener}.

\begin{lstlisting}[language=Java,caption={Inicio de Node.main(): config, dataset y bank}]
int port = (args.length > 0) ? Integer.parseInt(args[0]) : 8080;

NodeConfig config = NodeConfig.fromEnv();

Ledger ledger = new Ledger();
int loaded = Dataset.loadInto(ledger, Path.of(config.datasetPath()), config.idBase());
if (loaded <= 0) {
    throw new IllegalStateException("dataset loaded 0 accounts from "
            + config.datasetPath() + "; refusing to start with an empty ledger");
}
// ...
Bank bank = new Bank(ledger);
\end{lstlisting}

A continuacion el nodo se ramifica segun su rol, determinado por
\codigo{config.isLeader()}. Si es \textbf{lider} y hay Cloud Storage configurado
(\codigo{TES\_GCS\_KEYFILE} y \codigo{TES\_BUCKET}), construye un
\codigo{GcsStore} y un \codigo{Journal}, ejecuta la recuperacion en frio con
\codigo{journal.recover(bank::applyReplicated)} para reaplicar las
transferencias durables, registra el journal como \codigo{CommitListener} y
agrega un \emph{shutdown hook} que drena la cola durable al apagar. Tambien
levanta el \codigo{ReplicaFeed} en el puerto de replicacion y lo registra como
otro listener, de modo que cada commit en vivo se transmita por TCP a las
replicas conectadas. Si es \textbf{replica}, opcionalmente restaura su
\codigo{ReplicaCheckpoint} desde GCS \emph{antes} de iniciar la replicacion (para
que el \codigo{CATCHUP} arranque desde la marca exacta ya aplicada y no desde 0),
y luego un \codigo{ReplicaSync} se pone al dia desde el lider y aplica los
commits en vivo a traves del banco.

Por ultimo, \codigo{main()} arma la pila de seguridad
(\codigo{Tokens}, \codigo{CredentialStore}, \codigo{Authenticator}), construye el
proveedor de metricas \codigo{NodeStats}, registra la tabla de rutas con las
rutas exactas del contrato y arranca el servidor. La ultima linea bloquea el
hilo principal mientras los reactores y los workers atienden las peticiones.

\begin{lstlisting}[language=Java,caption={Tabla de rutas y arranque del servidor (Node.main())}]
Routes routes = new Routes();
routes.register("POST", "/api/register", new RegisterEndpoint(authenticator));
routes.register("POST", "/api/login", new LoginEndpoint(authenticator));
routes.register("GET", "/api/accounts/{id}", new BalanceEndpoint(authenticator, ledger));
routes.register("POST", "/api/transactions/transfer", new TransferEndpoint(authenticator, bank));
routes.register("GET", "/api/stats", new StatsEndpoint(stats));
routes.register("GET", "/panel", new PanelEndpoint(stats, config));
routes.register("GET", "/", new DashboardEndpoint());

Server server = new Server(port, routes, config.workers(), config.reactors());
server.start();

// Block the main thread while the IoLoop and workers serve requests.
Thread.currentThread().join();
\end{lstlisting}

El proveedor de metricas se instancia como
\codigo{NodeStats(ledger, bank, journal, config)} (ver \codigo{app/NodeStats.java}):
toma el numero de cuentas y el saldo total del \codigo{Ledger}, el conteo de
transferencias y el ultimo id del \codigo{Bank}, y el conteo en GCS del
\codigo{Journal} (que vale 0 en las replicas, donde el journal es null). Las
metricas de host (CPU, RAM, disco) se leen con Java crudo desde
\codigo{/proc/stat}, \codigo{/proc/meminfo} y el sistema de archivos, sin
bibliotecas externas.

\subsection{Configuracion por entorno (TES\_*)}

Toda la configuracion del nodo se inyecta por variables de entorno con prefijo
\codigo{TES\_}, leidas una sola vez en \codigo{NodeConfig.fromEnv()} dentro de
\codigo{app/NodeConfig.java}. El rol se deriva de \codigo{TES\_LEADER\_HOST}:
vacio o ausente significa que el proceso es el \textbf{lider} (nodo-1); cualquier
otro valor nombra al host del lider que esta \textbf{replica} debe seguir. El
secreto \codigo{TES\_JWT\_SECRET} debe ser el MISMO en los tres nodos para que un
token emitido en cualquiera valide en todos. La siguiente tabla resume cada
variable con su significado y su valor por defecto (los defaults vienen de las
constantes \codigo{DEFAULT\_*} de la clase).

\begin{center}
\begin{tabular}{|l|l|p{6.4cm}|}
\hline
\textbf{Variable} & \textbf{Default} & \textbf{Significado} \\
\hline
\codigo{TES\_DATASET} & (ninguno) & Ruta al CSV de cuentas que carga el nodo. \\
\codigo{TES\_JWT\_SECRET} & (ninguno) & Secreto HMAC para firmar/validar JWT; identico en los 3 nodos. \\
\codigo{TES\_NODE\_ID} & (ninguno) & Identificador legible del nodo, p.ej. \codigo{nodo-1}. \\
\codigo{TES\_PEERS} & \codigo{""} & URLs base de los pares, separadas por coma, para la agregacion del panel. \\
\codigo{TES\_LEADER\_HOST} & \codigo{""} & Host del lider a seguir; vacio = este nodo ES el lider. \\
\codigo{TES\_REPL\_PORT} & 9090 & Puerto TCP del feed de replicacion del lider. \\
\codigo{TES\_BUCKET} & (ninguno) & Bucket de Cloud Storage para el journal durable. \\
\codigo{TES\_GCS\_KEYFILE} & (ninguno) & Ruta al key file de la cuenta de servicio para autenticar contra GCS. \\
\codigo{TES\_WORKERS} & 16 & Tamano del pool de hilos que procesan peticiones. \\
\codigo{TES\_ID\_BASE} & 1 & Id asignado a la primera cuenta del dataset. \\
\codigo{TES\_REPLICA\_CHECKPOINT\_INTERVAL\_SECS} & 10 & Segundos entre flushes del checkpoint de la replica a GCS. \\
\codigo{TES\_REPLICA\_CHECKPOINT\_PATH} & (deriva) & Nombre del objeto del checkpoint; si es null usa \codigo{checkpoint/<nodeId>.json}. \\
\codigo{TES\_REACTORS} & 1 por CPU & Numero de reactores (selectores) del servidor HTTP. \\
\hline
\end{tabular}
\end{center}

La lectura usa dos ayudantes privados: \codigo{env(nombre, def)} devuelve el
default cuando la variable esta ausente o vacia, e \codigo{intEnv(nombre, def)}
parsea un entero y cae al default si la variable falta o no es numerica. El
metodo \codigo{isLeader()} centraliza la decision de rol: devuelve \codigo{true}
cuando \codigo{leaderHost} es null, vacio o solo espacios en blanco.

\begin{lstlisting}[language=Java,caption={Lectura de las variables TES\_* en NodeConfig.fromEnv()}]
public static NodeConfig fromEnv() {
    return new NodeConfig(
            env("TES_DATASET", null),
            env("TES_JWT_SECRET", null),
            env("TES_NODE_ID", null),
            env("TES_PEERS", ""),
            env("TES_LEADER_HOST", ""),
            intEnv("TES_REPL_PORT", DEFAULT_REPL_PORT),
            // ...
            intEnv("TES_WORKERS", DEFAULT_WORKERS),
            intEnv("TES_ID_BASE", DEFAULT_ID_BASE),
            // ...
            intEnv("TES_REACTORS", DEFAULT_REACTORS));
}
\end{lstlisting}

\subsection{Carga del dataset de cuentas}

Las cuentas se cargan desde el CSV con \codigo{Dataset.loadInto()} en
\codigo{core/Dataset.java}. El archivo de referencia es
\codigo{material-profesor/alumnos.csv}, entregado por el profesor y consumido tal
cual (no se regenera ni se edita). Es CSV plano de cuatro columnas y \emph{sin}
fila de encabezado, con el formato
\codigo{nombre,apellido1,apellido2,saldo}; por eso la primera linea tambien se
parsea, no se salta. Como no hay columna de id, el id de cada cuenta es su
posicion en el archivo contando desde \codigo{idBase} (normalmente 1, via
\codigo{TES\_ID\_BASE}): la primera linea de datos pasa a ser la cuenta
\codigo{idBase}, la segunda \codigo{idBase+1}, y asi sucesivamente. Que todos los
nodos lean el mismo archivo es justo lo que hace significativa la verificacion de
conservacion del dinero en el cluster.

El nombre del titular es la union con espacios de los tres campos de nombre. El
saldo se convierte de su texto decimal a centavos enteros con
\codigo{Money.toCents(String)}, de modo que ningun punto flotante toque jamas un
saldo: los importes se manejan internamente en centavos. Las lineas vacias o con
menos de cuatro columnas se descartan, y el metodo devuelve el numero de cuentas
efectivamente cargadas (el valor que \codigo{Node.main()} exige que sea mayor que
cero).

\begin{lstlisting}[language=Java,caption={Asignacion de id por indice de fila y saldo en centavos (Dataset.loadInto)}]
int id = idBase;
int loaded = 0;
// ...
String[] cols = line.split(",", -1);
if (cols.length < 4) {
    continue;
}
String owner = cols[0].trim() + " " + cols[1].trim() + " " + cols[2].trim();
long cents = Money.toCents(cols[3].trim());
ledger.add(new Account(id, owner, cents));
id++;
loaded++;
\end{lstlisting}

\section{Servidor HTTP en NIO puro}

El transporte HTTP de la tesoreria no se apoya en \codigo{com.sun.net.httpserver}
ni en ningun framework: es una pila propia escrita sobre \codigo{java.nio}, repartida
en ocho clases del paquete \codigo{net}. \codigo{net/Server.java} es la puerta de
entrada, \codigo{net/IoLoop.java} es el reactor que hace toda la E/S,
\codigo{net/Channel.java} guarda el estado de cada conexion, \codigo{net/RequestParser.java}
parsea HTTP de forma incremental, \codigo{net/Request.java} y \codigo{net/Reply.java}
modelan peticion y respuesta, y \codigo{net/Routes.java} con \codigo{net/Endpoint.java}
resuelven el ruteo. Este capitulo explica por que se construyo asi y como encajan las
piezas.

\subsection{Por que un servidor propio}

La razon principal es doble. Primero, la regla de "Java puro" del concurso prohibe
servidores de libreria; escribir el transporte sobre \codigo{java.nio} mantiene la
dependencia unicamente en el JDK. Segundo, y mas importante, el concurso evalua bajo
carga, y un servidor propio da control fino sobre la concurrencia: se decide cuantos
reactores corren, cuantos hilos de trabajo computan respuestas y como se reparte la
E/S entre nucleos.

Un servidor de libreria basado en un \codigo{ThreadPool} bloqueante (un hilo por
conexion) gasta un hilo del sistema operativo por cada socket abierto, lo que no
escala cuando hay miles de conexiones en keep-alive mayormente ociosas. El diseno NIO
multiplexa muchas conexiones sobre pocos hilos mediante un \codigo{Selector}: un solo
hilo reactor vigila cientos de sockets y solo despierta para los que tienen datos
listos. Asi se separa la E/S (barata, dirigida por eventos) del computo de la respuesta
(que puede tocar el banco y la replicacion), y cada parte se dimensiona por separado.

\subsection{Modelo de concurrencia: IoLoop y workers}

El diseno sigue el patron Reactor con pool de trabajadores (Doug Lea, patron 3) y usa
tres clases de hilo con responsabilidades disjuntas, lo que evita casi todo bloqueo
compartido.

El \codigo{Server} posee un \codigo{ServerSocketChannel}, un hilo aceptador dedicado,
un arreglo de reactores \codigo{IoLoop} (cada uno con su propio \codigo{Selector}) y un
pool compartido de hilos worker. El numero de reactores y de workers se fija en el
constructor; en produccion los provee \codigo{app/NodeConfig.java} desde las variables
de entorno \codigo{TES\_REACTORS} (por defecto uno por CPU) y \codigo{TES\_WORKERS}
(por defecto 16). Con \codigo{reactors == 1} se obtiene el reactor unico clasico, un
respaldo seguro.

El \textbf{hilo aceptador} corre \codigo{acceptLoop()}: se bloquea en \codigo{accept()}
y reparte cada conexion nueva entre los reactores en round-robin. Repartir aqui evita
que un solo reactor sature un nucleo mientras los demas estan ociosos, que es el cuello
de botella del reactor unico.

\begin{lstlisting}[language=Java,caption={El aceptador reparte cada conexion en round-robin (net/Server.java)}]
SocketChannel socket = serverChannel.accept();
if (socket == null) {
    continue;
}
socket.configureBlocking(false);
loops[next].register(socket);
next = (next + 1) % loops.length;
\end{lstlisting}

Cada \textbf{hilo reactor} (\codigo{IoLoop}) hace solo lectura y escritura para sus
conexiones; nunca calcula una \codigo{Reply} en linea. Su \codigo{run()} se bloquea en
\codigo{selector.select()} y, al despertar, primero drena dos colas:
\codigo{registerPending()} registra los sockets que el aceptador dejo en cola, y
\codigo{enableQueuedWrites()} activa \codigo{OP\_WRITE} para las llaves cuya respuesta
ya esta lista. Esto es deliberado: el registro de canales y las mutaciones del conjunto
de intereses (\codigo{interestOps}) ocurren unicamente en el hilo del reactor, de modo
que el aceptador y los workers se comunican mediante colas concurrentes
(\codigo{ConcurrentLinkedQueue}) y \codigo{selector.wakeup()}, manteniendo fuera del
diseno la clasica condicion de carrera entre hilos sobre el \codigo{Selector}.

Cuando \codigo{onRead()} lee bytes y el parser reporta una peticion completa, el reactor
no la atiende: pausa las lecturas de esa conexion con \codigo{key.interestOps(0)} y
envia el trabajo al pool con \codigo{workers.submit(...)}. Pausar las lecturas garantiza
una sola peticion en vuelo por conexion, lo que ordena las respuestas sin necesidad de
pipelining.

\begin{lstlisting}[language=Java,caption={El reactor delega el computo al worker (net/IoLoop.java)}]
boolean complete = channel.parser().feed(buffer);
buffer.compact();
if (complete) {
    Request request = channel.parser().take();
    channel.setProcessing(true);
    // Pause reads until this reply is flushed: one request in flight per
    // connection, which keeps responses ordered without pipelining.
    key.interestOps(0);
    workers.submit(() -> handle(key, channel, request));
}
\end{lstlisting}

El \textbf{hilo worker} ejecuta \codigo{handle()}: rutea con \codigo{routes.match(...)},
invoca al \codigo{Endpoint}, serializa la \codigo{Reply} a un \codigo{ByteBuffer} HTTP/1.1
y la encola en el \codigo{Channel}. Luego devuelve la llave al reactor anadiendola a
\codigo{writeReady} y llamando \codigo{selector.wakeup()}; el reactor activara
\codigo{OP\_WRITE} y \codigo{onWrite()} hara el flush real. El estado por conexion vive
en \codigo{Channel}, cuyo punto de encuentro es la cola \codigo{outbound}: el worker la
llena con \codigo{enqueue()} (sincronizado) y el reactor la vacia con
\codigo{peekOutbound()}/\codigo{popOutbound()} en cada evento de escritura. Si el buffer
de envio del socket se llena, \codigo{onWrite()} deja la escritura a medias y conserva
el interes en \codigo{OP\_WRITE} para reintentar despues.

\subsection{Parseo de peticiones y rutas}

Cada \codigo{Channel} tiene su propio \codigo{RequestParser}, un parser HTTP/1.1
incremental con maquina de estados explicita (\codigo{READ\_LINE}, \codigo{READ\_HEADERS},
\codigo{READ\_BODY}, \codigo{DONE}). Su metodo \codigo{feed(ByteBuffer)} consume los
bytes conforme llegan, sin re-escanear el buffer desde cero, y devuelve \codigo{true}
cuando hay una peticion completa disponible via \codigo{take()}. En \codigo{READ\_LINE}
y \codigo{READ\_HEADERS} acumula una linea terminada en CRLF; al llegar el \verb|\n|,
\codigo{parseRequestLine()} divide "METHOD destino HTTP/x.y" (separando el destino en
\codigo{path} y \codigo{query} por el primer \codigo{?}) y \codigo{parseHeader()} guarda
cada cabecera con su nombre en minusculas. La linea vacia cierra las cabeceras: si hay
\codigo{Content-Length} se pasa a \codigo{READ\_BODY} y se leen exactamente esos bytes;
si no, la peticion ya esta completa. \codigo{take()} entrega un \codigo{Request} inmutable
y resetea la maquina para la siguiente peticion sobre la misma conexion keep-alive. El
parser tambien decide el keep-alive (\codigo{computeKeepAlive()}): respeta la cabecera
\codigo{Connection} y, en su ausencia, mantiene viva la conexion solo en HTTP/1.1.

El \codigo{Request} resultante es una vista inmutable: expone \codigo{method()},
\codigo{path()}, \codigo{query()}, \codigo{header(name)} (insensible a mayusculas),
\codigo{body()}/\codigo{bodyText()} y \codigo{pathParam(name)}. La firma del handler la
define \codigo{net/Endpoint.java}, una interfaz funcional con un solo metodo:

\begin{lstlisting}[language=Java,caption={Contrato de un manejador de peticiones (net/Endpoint.java)}]
@FunctionalInterface
public interface Endpoint {

    /** Handles one request and produces its reply. */
    Reply serve(Request req);
}
\end{lstlisting}

El ruteo lo hace \codigo{Routes}, un router por metodo y path. Cada patron se parte en
segmentos por \codigo{/}; un segmento escrito como \codigo{\{name\}} es una variable de
un solo segmento y los demas deben coincidir por igualdad exacta de cadena (es coincidencia
exacta de segmentos, no de prefijo mas largo). \codigo{register(method, pattern, endpoint)}
compila el patron y \codigo{match(method, path)} lo resuelve, capturando las variables.
Las semanticas de resolucion son precisas: si coinciden el metodo y la forma de segmentos
devuelve un \codigo{Match} con estado 200 y el \codigo{Endpoint}; si algun patron coincide
con el path pero ningun metodo encaja devuelve 405; y si ningun patron coincide, 404.

\begin{lstlisting}[language=Java,caption={Coincidencia de segmentos y captura de variables (net/Routes.java)}]
for (int i = 0; i < segments.length; i++) {
    String pat = route.segments[i];
    if (isVariable(pat)) {
        params.put(pat.substring(1, pat.length() - 1), segments[i]);
    } else if (!pat.equals(segments[i])) {
        shapeMatches = false;
        break;
    }
}
\end{lstlisting}

En \codigo{IoLoop.handle()}, cuando \codigo{match.status} es 200 se llama
\codigo{request.bindPathParams(match.params)} para inyectar las variables capturadas y
luego \codigo{match.endpoint.serve(request)}. Un estado distinto de 200 se traduce a
\codigo{Reply.status(match.status)} (404 o 405); si el endpoint devuelve \codigo{null} o
lanza una \codigo{RuntimeException}, el reactor responde 500.

\subsection{Construccion de respuestas}

La respuesta se modela con \codigo{net/Reply.java}, un objeto que el worker construye y
el reactor serializa despues. Lleva el codigo de estado, el \codigo{Content-Type} y el
cuerpo en bytes. Sus factorias estaticas cubren los tipos que emite el servicio:
\codigo{Reply.json(status, json)} fija \codigo{application/json; charset=utf-8} y
codifica el cuerpo en UTF-8; \codigo{Reply.text(...)} produce \codigo{text/plain};
\codigo{Reply.html(...)} sirve \codigo{text/html} a partir de bytes crudos (por ejemplo
una pagina leida de recursos); y \codigo{Reply.status(int)} arma una respuesta solo de
estado con cuerpo vacio, la que usa el reactor para 404, 405 y 500.

\begin{lstlisting}[language=Java,caption={Factoria JSON de respuestas (net/Reply.java)}]
public static Reply json(int status, String json) {
    return new Reply(status, "application/json; charset=utf-8",
            json.getBytes(StandardCharsets.UTF_8));
}
\end{lstlisting}

\codigo{Reply} ofrece ademas \codigo{close()}, que marca la conexion para cerrarse tras
enviar la respuesta (encadenable porque devuelve \codigo{this}). La escritura efectiva
del HTTP no ocurre en \codigo{Reply}: el reactor la serializa en \codigo{IoLoop.serialize()},
que arma la linea de estado \codigo{HTTP/1.1 <codigo> <razon>} con las cabeceras
\codigo{Content-Type}, \codigo{Content-Length} y \codigo{Connection} (\codigo{keep-alive}
o \codigo{close}, segun lo decidido entre \codigo{request.wantsClose()} y
\codigo{reply.isClose()}), seguidas del cuerpo. La razon textual la traduce
\codigo{reason(int)} para los codigos que el servicio emite (200, 201, 400, 401, 403,
404, 405, 409, 500). De este modo \codigo{Reply} permanece como un valor inmutable y
agnostico al cable, y toda la mecanica de bytes y keep-alive queda concentrada en el
reactor.

\section{Nucleo del banco: dinero, ledger e invariante}

El paquete \codigo{core} concentra el dominio del mini-banco: como se representa
el dinero, donde viven los saldos, como se mueve dinero de una cuenta a otra sin
condiciones de carrera y que invariante debe respetarse siempre. Es el corazon
de la rubrica de Consistencia (25 puntos): el saldo global debe conservarse, sin
inconsistencias, aun bajo carga concurrente. Todo lo que sigue se apoya en una
decision de diseno simple pero estricta: el dinero nunca se representa con punto
flotante.

\subsection{Dinero exacto en centavos}

Un saldo jamas se guarda como \codigo{double}: se guarda como un entero de
centavos (\codigo{long}). Asi se evitan los errores de redondeo binario tipicos
del punto flotante, donde una cantidad como 0.10 no tiene representacion exacta y
las sumas acumulan deriva. Al trabajar en centavos enteros, toda operacion
aritmetica (sumar, restar, comparar) es exacta y reproducible en cualquier nodo.

Por eso los saldos viven como \codigo{long} en \codigo{core/Account.java}
(campo \codigo{balanceCents}) y los montos de cada transferencia tambien, en el
campo \codigo{cents} de \codigo{core/Transfer.java}. La clase
\codigo{core/Money.java} no almacena dinero ni hace aritmetica de saldos: es un
ayudante estatico (constructor privado, sin instancias) que centraliza la unica
conversion entre la cadena decimal que ve el humano (con dos decimales) y la
representacion interna en centavos. Se extrajo a proposito para que los
endpoints, el cargador de datos y el codec del cluster nunca incrusten
aritmetica de \codigo{BigDecimal} ad hoc.

La conversion de texto a centavos usa \codigo{BigDecimal} con redondeo
\codigo{HALF\_UP}, de modo que el parseo nunca introduce error de coma flotante:

\begin{lstlisting}[language=Java,caption={Conversion exacta de decimal a centavos en Money.toCents}]
public static long toCents(String decimal) {
    return new BigDecimal(decimal)
            .movePointRight(2)
            .setScale(0, RoundingMode.HALF_UP)
            .longValueExact();
}
\end{lstlisting}

El camino inverso, \codigo{Money.toDecimal(long)}, reconstruye la cadena con
exactamente dos decimales usando \codigo{movePointLeft(2)} y
\codigo{setScale(2)}. La aritmetica real de mover dinero (la resta del origen y
la suma al destino) no ocurre en \codigo{Money}, sino en el ledger, sobre esos
mismos \codigo{long} de centavos.

\subsection{Cuentas y ledger}

Una \codigo{core/Account.java} es una cuenta del banco: su id y su titular
(\codigo{owner}) son inmutables, y el unico estado mutable es el saldo en
centavos. Ese saldo solo puede leerse o escribirse manteniendo el monitor
intrinseco de la propia cuenta; por eso tanto el lector \codigo{balanceCents()}
como el escritor \codigo{setBalanceCents(long)} estan declarados
\codigo{synchronized}. No hay un objeto candado por campo: el ledger toma los
monitores intrinsecos de las dos cuentas que participan en una transferencia,
ordenados por id, para mantener el protocolo libre de deadlock.

El \codigo{core/Ledger.java} guarda el conjunto de cuentas y realiza la mecanica
cruda de mover dinero entre dos de ellas; deliberadamente no sabe nada de
secuencias, log de commits ni replicacion (eso vive en \codigo{Bank}). Las
cuentas se almacenan en un \codigo{ConcurrentHashMap<Integer, Account>}, de modo
que las lecturas (\codigo{get}) y el alta de cuentas (\codigo{add}) son libres de
candado. La consulta de un saldo individual es simplemente
\codigo{ledger.get(id).balanceCents()}. La suma global se obtiene con
\codigo{totalCents()}, que recorre todas las cuentas acumulando sus saldos como
una lectura consistente: aun bajo movimientos concurrentes ese total es
exactamente el valor cargado al inicio, porque el escaneo toma un snapshot
atomico (lo detalla la subseccion \emph{Snapshot atomico del total}). Observarlo
bajo carga es la forma en que el sistema verifica que las transferencias
conservan el dinero.

Todos los nodos parten del mismo estado gracias a \codigo{core/Dataset.java},
que carga el dataset compartido de cuentas en un ledger. El archivo es un CSV
plano de cuatro columnas sin encabezado
(\codigo{firstName,lastName1,lastName2,balance}) y sin columna de id: el id de
cada cuenta es su posicion en el archivo contando desde \codigo{idBase}
(normalmente 1), asi que la primera linea de datos se parsea, no se omite. El
titular se forma uniendo los tres campos de nombre con espacios, y el saldo se
convierte de su texto decimal a centavos enteros con \codigo{Money.toCents}, de
modo que ningun punto flotante toca jamas un saldo. Como cada nodo lee el mismo
archivo, todos arrancan con saldos identicos, y por eso la comprobacion de
conservacion de dinero a lo largo del cluster es significativa.

\subsection{Transferencia atomica}

La orquestacion de una transferencia solicitada por el cliente vive en
\codigo{Bank.transfer(int, int, long)} de \codigo{core/Bank.java}. El banco
primero aplica las validaciones que son politica del banco (no mecanica del
ledger): rechaza una transferencia a la misma cuenta y un monto no positivo.
Solo si esas comprobaciones pasan, delega el debito y el credito al ledger y, en
el mismo paso atomico (dentro del candado del movimiento), le asigna el siguiente
numero de secuencia; luego la registra en el log y la entrega a los listeners:

\begin{lstlisting}[language=Java,caption={Bank.transfer: politica, movimiento delegado y commit}]
public long transfer(int from, int to, long cents) throws TransferException {
    if (from == to) {
        throw TransferException.selfTransfer();
    }
    if (cents <= 0L) {
        throw TransferException.badAmount();
    }
    // La secuencia se asigna DENTRO del candado del movimiento (atomica con el).
    long seq = ledger.move(from, to, cents, sequence::incrementAndGet);
    Transfer committed = new Transfer(seq, from, to, cents);
    // ...
}
\end{lstlisting}

La parte critica para la consistencia, la atomicidad y la ausencia de deadlock,
no esta en \codigo{Bank} sino en el metodo privado \codigo{apply} del ledger, al
que llega \codigo{ledger.move}. Ese metodo busca ambas cuentas y, para mover el
dinero como un solo paso atomico, toma los dos monitores en un orden fijo
determinado por el id (siempre el menor primero). Ese orden global es lo que
impide el deadlock: si dos transferencias inversas tocan el mismo par de cuentas
en sentidos opuestos, ambas intentan adquirir los candados en la misma
secuencia, de modo que nunca se bloquean mutuamente esperando candados cruzados.

\begin{lstlisting}[language=Java,caption={Ledger.apply: orden de candados por id y movimiento atomico}]
Account first = (from < to) ? src : dst;
Account second = (from < to) ? dst : src;
synchronized (first) {
    synchronized (second) {
        if (enforceFunds && src.balanceCents() < cents) {
            throw TransferException.lowBalance(from);
        }
        src.setBalanceCents(src.balanceCents() - cents);
        dst.setBalanceCents(dst.balanceCents() + cents);
        // Estampa la secuencia aqui dentro: el orden de commit coincide
        // con el orden en que se movio el dinero.
        return (sequencer != null) ? sequencer.getAsLong() : 0L;
    }
}
\end{lstlisting}

Como ambos monitores se mantienen durante todo el movimiento, el par de saldos
cambia junto, sin que otro hilo observe un estado intermedio donde el dinero ya
salio del origen pero aun no entro al destino. El rechazo nunca se devuelve como
valor de retorno, sino como una \codigo{core/TransferException.java}, una
excepcion verificada con un \codigo{code()} de un vocabulario pequeno y cerrado:
\codigo{self\_transfer}, \codigo{bad\_amount}, \codigo{no\_such\_account} y
\codigo{low\_balance}. Asi un llamador no puede continuar por accidente despues
de un movimiento rechazado, y la capa HTTP mapea esos codigos a estados y a un
cuerpo JSON de error. La \codigo{core/Transfer.java} resultante es un
\codigo{record} inmutable (\codigo{seq}, \codigo{from}, \codigo{to},
\codigo{cents}): es la unidad de replicacion y durabilidad que el codec del
cluster y el journal en GCS serializan.

\subsection{Snapshot atomico del total (lectura consistente)}

Sumar las 820k cuentas mientras el banco recibe transferencias no es inocuo: un
escaneo ingenuo puede caer justo entre el debito y el credito de una misma
transferencia (origen ya restado, destino aun no sumado) y reportar un total
\emph{roto} que no corresponde a ningun estado real. Para que
\codigo{Ledger.totalCents()} sea siempre una lectura consistente, el ledger anade
un segundo mecanismo de sincronizacion, separado de los monitores de cuenta: un
\codigo{ReentrantReadWriteLock} llamado \codigo{snapshotLock}.

La clave es una inversion deliberada de los lados del candado. Cada movimiento de
dinero (los caminos \codigo{move}, \codigo{settle} y el \codigo{apply} privado
que comparten) toma el lado \textbf{COMPARTIDO} (read) de \codigo{snapshotLock}
mientras hace su debito y su credito; el escaneo de \codigo{totalCents()} toma el
lado \textbf{EXCLUSIVO} (write). Es decir, los lectores-del-candado son los
escritores-del-dinero, y el escritor-del-candado es el unico lector global. Como
el lado compartido admite cualquier numero de titulares simultaneos, todas las
transferencias siguen ejecutandose en paralelo entre si; pero el lado exclusivo
excluye a todas, de modo que la suma no puede solaparse jamas con una
transferencia a medio aplicar. El escaneo observa, por construccion, un estado
donde cada transferencia esta completa o aun no empezada: un \emph{snapshot}
atomico.

\begin{lstlisting}[language=Java,caption={Ledger: lado compartido en el movimiento, lado exclusivo en el escaneo}]
// En apply(...): el debito+credito corre bajo el lado COMPARTIDO,
// asi muchos transferencias avanzan a la vez. El readLock se toma ANTES
// de los monitores de cuenta (el monitor de cuenta es el sumidero del grafo).
snapshotLock.readLock().lock();
try {
    synchronized (first) {
        synchronized (second) {
            src.setBalanceCents(src.balanceCents() - cents);
            dst.setBalanceCents(dst.balanceCents() + cents);
            return (sequencer != null) ? sequencer.getAsLong() : 0L;
        }
    }
} finally {
    snapshotLock.readLock().unlock();
}

// En totalCents(): el escaneo corre bajo el lado EXCLUSIVO, que excluye
// a todo movimiento, por lo que nunca ve un debito sin su credito.
public long totalCents() {
    snapshotLock.writeLock().lock();
    try {
        long sum = 0L;
        for (Account a : accounts.values()) {
            sum += a.balanceCents();
        }
        return sum;
    } finally {
        snapshotLock.writeLock().unlock();
    }
}
\end{lstlisting}

Las dos capas de candados componen sin riesgo de deadlock porque forman un grafo
dirigido aciclico: el \codigo{readLock} se adquiere \emph{antes} que los monitores
de cuenta, y el monitor de cuenta es siempre el sumidero (nunca se toma el
\codigo{snapshotLock} teniendo ya un monitor de cuenta). Asi, el orden por id que
ya evitaba el deadlock entre transferencias inversas se mantiene intacto, y el
nuevo candado solo agrega un nivel por encima.

\textbf{Por que el invariante queda exacto en los tres nodos.} Como toda
transferencia es suma-cero (un debito acompanado de un credito igual), el total
verdadero es invariante: no cambia jamas. Y como el escaneo de cada nodo corre
bajo el lado exclusivo, no puede capturar una transferencia a medio aplicar; por
tanto reporta exactamente ese total invariante, no un valor transitorio. Los tres
nodos parten del mismo dataset, de modo que los tres escaneos atomicos devuelven
la misma constante aun bajo carga. Eso es justo lo que el sello de consistencia
del dashboard compara: con el snapshot atomico, el sello permanece verde bajo
trafico sostenido en lugar de parpadear por lecturas rotas espurias.

\textbf{Por que \codigo{fair} evita la inanicion del escaneo.} El
\codigo{snapshotLock} se construye en modo justo (\codigo{new
ReentrantReadWriteLock(true)}). Sin justicia, un flujo continuo de transferencias
(lectores-del-candado) podria mantener el lado compartido siempre tomado y
\emph{matar de hambre} al escaneo (escritor-del-candado), que nunca conseguiria
el lado exclusivo. Con justicia, el escritor que espera se atiende en orden de
llegada, asi que el endpoint \codigo{/panel} no puede colgarse bajo carga
sostenida de transferencias.

\subsection{Invariante de saldo y numero de secuencia}

El invariante central es la conservacion del dinero: como cada movimiento es un
debito en una cuenta acompanado de un credito por la misma cantidad en otra, la
suma global de saldos permanece constante. Ese es precisamente el valor que
expone \codigo{Ledger.totalCents()}, leido como un snapshot atomico, y su
constancia bajo carga es lo que demuestra que las transferencias no crean ni
destruyen dinero. El ledger se
mantiene a proposito estrecho (solo aplica debito y credito equilibrados) para
que ese invariante sea facil de razonar.

El numero de secuencia es la otra mitad de la consistencia, y no vive en el
ledger sino en \codigo{core/Bank.java}, como un \codigo{AtomicLong sequence}. Es
monotono y solo avanza en un commit real: la secuencia
(\codigo{sequence.incrementAndGet()}) se asigna \emph{dentro} del candado del
movimiento, atomica con el debito y el credito, de modo que una transferencia
rechazada nunca consume un numero y, ademas, el orden de las secuencias coincide
con el orden en que el dinero se movio. Eso importa para las replicas: como sirven
lecturas mientras reproducen en orden de secuencia, estampar la secuencia junto al
movimiento evita que una cuenta se vea \emph{transitoriamente} en negativo (un
debito reproducido antes del credito que lo fondea). Esa secuencia, estampada en
cada \codigo{Transfer}, es la base de la replicacion y de la verificacion de
consistencia: ordena los commits y permite a una replica saber hasta donde esta al
dia. En el camino de replica,
\codigo{Bank.applyReplicated(Transfer)} no genera secuencia nueva sino que la
recibe con la entrada y eleva los contadores con
\codigo{accumulateAndGet(t.seq(), Math::max)}, de modo que reaplicar una entrada
duplicada nunca retrocede el avance.

Cada commit aceptado se registra en el \codigo{core/TransferLog.java}, un
historial append-only y thread-safe. Conserva la secuencia ordenada de commits
para que la alimentacion de replicas pueda reproducir el historial de catch-up a
una replica recien conectada. Las altas ocurren en el camino de commit a la tasa
plena de transferencias (20\% del trafico del concurso), asi que la lista interna
es un \codigo{ArrayList} protegido por el monitor del objeto: anexar es O(1)
amortizado. Una lista copy-on-write seria un error aqui porque copiaria todo el
arreglo en cada anexion, volviendo la carga sostenida un costo de copia
cuadratico. Los escaneos de catch-up (\codigo{since}) son raros (solo cuando una
replica se conecta) y devuelven una copia fresca de la cola que coincide, de modo
que el llamador la itera sin sostener el candado.

Por ultimo, \codigo{core/CommitListener.java} es la interfaz funcional que
notifica un commit. Los listeners se registran con \codigo{Bank.addCommitListener}
y, en el codigo, el fan-out ocurre unicamente en el camino de cliente del lider:
al final de \codigo{Bank.transfer} se recorre la lista (un
\codigo{CopyOnWriteArrayList}, para no bloquear el registro) y se invoca
\codigo{onCommit(committed)} en cada uno, en orden de registro. Los listeners
concretos son el journal en GCS (durabilidad) y la alimentacion de replicas
(replicacion en vivo). La via \codigo{applyReplicated} no vuelve a notificar a
los listeners, porque ya esta consumiendo el flujo del lider y no debe
re-emitirlo.

\section{Seguridad: JWT y credenciales}

La autenticación del API se apoya en tres piezas que viven en el paquete
\codigo{security/}: un emisor/verificador de JWT (\codigo{security/Tokens.java}),
un hashing de contraseñas (\codigo{security/Passwords.java}) con su almacén de
credenciales (\codigo{security/Credential.java} y
\codigo{security/CredentialStore.java}), y un orquestador
(\codigo{security/Authenticator.java}) que une todo y protege los endpoints. El
JWT es real (firmado con HMAC y verificado en cada petición protegida), no
decorativo: el login valida la credencial y el registro guarda el hash, nunca la
contraseña en claro.

\subsection{Emision y verificacion de JWT}

\codigo{security/Tokens.java} es un servicio de instancia que envuelve la
librería \codigo{com.auth0.jwt} (java-jwt). Se construye con el secreto HMAC
compartido y, a partir de él, crea un \codigo{Algorithm.HMAC256(secret)} y un
\codigo{JWTVerifier} ya fijado al emisor \codigo{tesoreria}. Construir el
verificador una sola vez en el constructor evita recompilar las reglas de
verificación en cada petición.

\begin{lstlisting}[language=Java,caption={Tokens: constructor, emision y validacion}]
public Tokens(String secret) {
    this.algorithm = Algorithm.HMAC256(secret);
    this.verifier = JWT.require(algorithm).withIssuer(ISSUER).build();
}

public String issue(String subject) {
    return JWT.create()
            .withIssuer(ISSUER)
            .withSubject(subject)
            .sign(algorithm);
}

public String validate(String token) {
    DecodedJWT decoded = verifier.verify(token);
    return decoded.getSubject();
}
\end{lstlisting}

\textbf{Qué claims lleva.} El método \codigo{issue(subject)} firma un token
compacto con dos claims: el emisor (\codigo{iss = tesoreria}) y el sujeto
(\codigo{sub}), que es el nombre de usuario autenticado. El token no incluye un
claim de expiración (\codigo{exp}): no se llama a \codigo{withExpiresAt}, de modo
que el token es válido mientras lo sea la firma. La firma es HMAC-SHA256 sobre el
secreto compartido, lo que garantiza integridad y autenticidad pero no
confidencialidad (el payload va en Base64, no cifrado), por lo que no debe
guardarse información sensible en los claims.

\textbf{Cómo se verifica.} En cada petición protegida, \codigo{validate(token)}
llama a \codigo{verifier.verify(token)}, que comprueba la firma HMAC256 y que el
emisor sea \codigo{tesoreria}; si algo no cuadra, java-jwt lanza una
\codigo{JWTVerificationException} (subtipo de \codigo{RuntimeException}). Si la
verificación pasa, devuelve el sujeto, es decir, la identidad que viajaba en el
token.

\textbf{Secreto compartido entre nodos.} El secreto se inyecta desde fuera: en
\codigo{app/Node.java} se construye \codigo{new Tokens(config.jwtSecret())}, y ese
valor proviene de la variable de entorno \codigo{TES\_JWT\_SECRET} leída por
\codigo{app/NodeConfig.java}. Como documenta ese archivo, \codigo{TES\_JWT\_SECRET}
debe tener \emph{el mismo valor en los tres nodos}, de forma que un token emitido
en cualquier nodo (por ejemplo, tras hacer login en el líder) valide en todos. Es
la propiedad clave de un secreto HMAC simétrico: quien conoce el secreto puede
tanto firmar como verificar.

\subsection{Contrasenas y credenciales}

Las contraseñas nunca se almacenan en texto plano. \codigo{security/Passwords.java}
es un ayudante sin estado, con dos métodos estáticos respaldados por jBCrypt
(\codigo{org.mindrot.jbcrypt.BCrypt}), que esconde la generación de la sal y la
comparación tras una API mínima.

\begin{lstlisting}[language=Java,caption={Passwords: hash y verificacion con BCrypt}]
public static String encode(String raw) {
    return BCrypt.hashpw(raw, BCrypt.gensalt());
}

public static boolean matches(String raw, String hash) {
    return BCrypt.checkpw(raw, hash);
}
\end{lstlisting}

\codigo{encode(raw)} genera una sal nueva con \codigo{BCrypt.gensalt()} y produce
el hash BCrypt, que ya lleva la sal incrustada en su cadena. Se usa BCrypt en
lugar de un hash simple porque es deliberadamente lento (factor de coste) y salado
por diseño, lo que dificulta los ataques de fuerza bruta y de tablas
precalculadas. \codigo{matches(raw, hash)} re-deriva el hash del candidato con la
misma sal del hash guardado y compara; \codigo{BCrypt.checkpw} encapsula esa
comparación. Como cada llamada a \codigo{encode} usa una sal distinta, dos
usuarios con la misma contraseña obtienen hashes diferentes.

Una credencial almacenada se modela con \codigo{security/Credential.java}, un
\codigo{record} inmutable que empareja el \codigo{username} con su
\codigo{passwordHash} BCrypt y, por contrato, nunca contiene la contraseña en
claro:

\begin{lstlisting}[language=Java,caption={Credential: record inmutable}]
public record Credential(String username, String passwordHash) {
}
\end{lstlisting}

El registro de credenciales vive en \codigo{security/CredentialStore.java}, un
almacén en memoria seguro para hilos basado en un
\codigo{ConcurrentHashMap<String, Credential>} indexado por nombre de usuario. La
operación de alta es atómica: \codigo{register(username, passwordHash)} usa
\codigo{putIfAbsent}, de modo que un nombre de usuario sólo puede reclamarse una
vez aunque dos peticiones de registro lleguen en paralelo.

\begin{lstlisting}[language=Java,caption={CredentialStore: alta atomica y busqueda}]
public boolean register(String username, String passwordHash) {
    Credential created = new Credential(username, passwordHash);
    return byUsername.putIfAbsent(username, created) == null;
}

public Credential find(String username) {
    return byUsername.get(username);
}
\end{lstlisting}

\codigo{register} devuelve \codigo{true} si el alta tuvo éxito y \codigo{false} si
el usuario ya existía (cuando \codigo{putIfAbsent} no es \codigo{null}); y
\codigo{find(username)} devuelve la credencial guardada o \codigo{null} si no hay
ninguna. Al ser un almacén en memoria, las credenciales son locales al proceso y
no se replican entre nodos junto con el libro mayor.

\subsection{Autenticador y proteccion de endpoints}

\codigo{security/Authenticator.java} es la fachada que coordina alta de cuentas,
login y autorización de peticiones. Se construye inyectándole un
\codigo{CredentialStore} (quién existe) y un \codigo{Tokens} (prueba de identidad);
en \codigo{app/Node.java} se crea con \codigo{new Authenticator(credentials,
tokens)}.

\textbf{Alta y login.} \codigo{signup(username, password)} hashea la contraseña
con \codigo{Passwords.encode} y delega en \codigo{store.register}, devolviendo
\codigo{false} si el usuario estaba tomado. \codigo{login(username, password)}
busca la credencial, verifica la contraseña con \codigo{Passwords.matches} y, sólo
si todo concuerda, emite un JWT con \codigo{tokens.issue(username)}; en cualquier
fallo (usuario inexistente o contraseña incorrecta) devuelve \codigo{null}, sin
distinguir el caso para no filtrar qué usuarios existen.

\textbf{Autorización por Bearer.} \codigo{authorize(req)} extrae la cabecera
\codigo{Authorization}, exige el esquema \codigo{Bearer} y valida el token:

\begin{lstlisting}[language=Java,caption={Authenticator.authorize: validacion del Bearer}]
public String authorize(Request req) {
    String header = req.header("Authorization");
    if (header == null) {
        return null;
    }
    String[] parts = WHITESPACE.split(header.trim(), 2);
    if (parts.length < 2 || !SCHEME.equalsIgnoreCase(parts[0])) {
        return null;
    }
    String token = parts[1].trim();
    if (token.isEmpty()) {
        return null;
    }
    try {
        return tokens.validate(token);
    } catch (RuntimeException e) {
        return null;
    }
}
\end{lstlisting}

El método es tolerante con el formato según el RFC 7235: el esquema se compara con
\codigo{equalsIgnoreCase} (acepta \codigo{Bearer}, \codigo{bearer} o
\codigo{BEARER}) y los espacios sobrantes se absorben partiendo con un patrón de
espacios precompilado (\codigo{WHITESPACE}). Devuelve el sujeto del token cuando
la firma y el emisor son válidos, y \codigo{null} en todos los casos de fallo:
cabecera ausente, esquema equivocado, token vacío o cualquier
\codigo{JWTVerificationException} capturada como \codigo{RuntimeException}.

\textbf{Rutas públicas frente a protegidas.} El registro de rutas se hace en
\codigo{app/Node.java}. Son \emph{públicas} \codigo{POST /api/register} (atendida
por \codigo{RegisterEndpoint}) y \codigo{POST /api/login} (atendida por
\codigo{LoginEndpoint}): no requieren token porque su propósito es precisamente
obtenerlo. Son \emph{protegidas} las que operan sobre cuentas, como
\codigo{GET /api/accounts/\{id\}} (\codigo{BalanceEndpoint}) y
\codigo{POST /api/transactions/transfer} (\codigo{TransferEndpoint}). Cada una
empieza llamando a \codigo{auth.authorize(req)} y, si devuelve \codigo{null},
corta con \codigo{Reply.status(401)} antes de tocar la lógica de negocio:

\begin{lstlisting}[language=Java,caption={BalanceEndpoint: corte 401 si el token falla}]
if (auth.authorize(req) == null) {
    return Reply.status(401);
}
\end{lstlisting}

Así, una petición sin cabecera \codigo{Authorization}, con un esquema distinto de
Bearer o con un token mal firmado o de otro emisor recibe un \codigo{401
Unauthorized}, mientras que un token válido deja pasar la petición a la operación
correspondiente.

\section{API REST y endpoints}

Este capítulo documenta el contrato REST del nodo: cada endpoint con su método,
ruta, cuerpo, respuesta y códigos de estado. Toda la API se sirve sobre el mismo
router (\codigo{net/Routes.java}) y dos de los endpoints exigen un JWT,
siguiendo la estructura REST + JWT que pide la rúbrica.

\subsection{Mapa de rutas}

El enrutador vive en \codigo{net/Routes.java}. Es un \emph{router} por método y
ruta: cada patrón se parte en segmentos separados por \codigo{/} y un segmento
escrito como \codigo{\{nombre\}} es una variable de un solo segmento; el resto
debe coincidir por igualdad exacta de cadena (no es coincidencia por prefijo más
largo). El método \codigo{Routes.match()} resuelve un método y una ruta concreta
contra las rutas registradas con una semántica clara de tres resultados: si el
método y la forma de los segmentos coinciden devuelve el \codigo{Endpoint} con
estado \codigo{200}; si algún patrón coincide con la ruta pero ningún método
coincide devuelve \codigo{405} (\emph{method not allowed}); y si ningún patrón
coincide con la ruta devuelve \codigo{404}.

El cableado real de rutas a \emph{endpoints} ocurre en \codigo{app/Node.java},
mediante llamadas a \codigo{routes.register(...)}. La siguiente tabla resume cada
endpoint registrado, con su método, ruta, si exige autenticación con \emph{token}
y su propósito.

\begin{center}
\begin{tabular}{l l c l}
\textbf{Método} & \textbf{Ruta} & \textbf{Auth} & \textbf{Propósito} \\
\hline
POST & \codigo{/api/register} & No & Crea una credencial de usuario. \\
POST & \codigo{/api/login} & No & Verifica credencial y emite un JWT. \\
GET & \codigo{/api/accounts/\{id\}} & Sí & Devuelve el saldo de una cuenta. \\
POST & \codigo{/api/transactions/transfer} & Sí & Transfiere dinero entre cuentas. \\
GET & \codigo{/api/stats} & No & Métricas de \emph{este} nodo. \\
GET & \codigo{/panel} & No & Métricas agregadas del clúster. \\
GET & \codigo{/} & No & Sirve el tablero de una sola página. \\
\end{tabular}
\end{center}

Solo los dos endpoints de negocio (\codigo{/api/accounts/\{id\}} y
\codigo{/api/transactions/transfer}) exigen un \emph{token} válido; los de
autenticación y los de monitoreo son públicos. La columna \emph{Auth} se valida
con \codigo{Authenticator.authorize()} dentro de cada endpoint, no en el router.

\subsection{Autenticacion: register y login}

\subsubsection*{Registro de usuario}

\codigo{api/RegisterEndpoint.java} atiende \codigo{POST /api/register} y crea una
credencial nueva. Acepta un cuerpo JSON con la forma
\codigo{\{"username":"...","password":"..."\}}, lo parsea con \codigo{Gson} y
delega la persistencia en \codigo{Authenticator.signup()}, que cifra la
contraseña y la guarda en el \codigo{CredentialStore}. El endpoint primero
verifica el método (responde \codigo{405} si no es \codigo{POST}); si el JSON está
mal formado o falta \codigo{username}/\codigo{password} responde \codigo{400} con
\codigo{\{"error":"bad\_request"\}}; si el usuario aún no existía devuelve
\codigo{201} con \codigo{\{"status":"created"\}}; y si el nombre ya estaba tomado
devuelve \codigo{409} con \codigo{\{"error":"username\_taken"\}}. La distinción
entre creado (\codigo{201}) y duplicado (\codigo{409}) la decide directamente el
valor booleano que regresa \codigo{auth.signup(username, password)}.

\subsubsection*{Inicio de sesión}

\codigo{api/LoginEndpoint.java} atiende \codigo{POST /api/login}: verifica la
credencial y, en caso de éxito, emite un JWT. Recibe el mismo cuerpo
\codigo{\{"username":"...","password":"..."\}} y delega en
\codigo{Authenticator.login()}, que busca la credencial en el
\codigo{CredentialStore}, compara la contraseña con \codigo{Passwords.matches()} y
solo si coinciden mina un \emph{token} firmado con \codigo{Tokens.issue()}. Es
decir, el login valida credenciales reales: si el usuario no existe o la
contraseña no coincide, \codigo{login()} regresa \codigo{null} y el endpoint
responde \codigo{401} con \codigo{\{"error":"invalid\_credentials"\}}. Con
credenciales válidas responde \codigo{200} con \codigo{\{"token":"<jwt>"\}}; un
JSON mal formado produce \codigo{400} y un método distinto de \codigo{POST}
produce \codigo{405}. El cliente reutiliza ese \emph{token} en la cabecera
\codigo{Authorization: Bearer <jwt>} para las operaciones protegidas.

\begin{lstlisting}[language=none,caption={Contrato JSON de POST /api/login (peticion y respuesta)}]
POST /api/login
Content-Type: application/json

{"username":"juan","password":"secreto"}

200 OK
{"token":"<jwt>"}

401 Unauthorized
{"error":"invalid_credentials"}
\end{lstlisting}

\subsection{Operaciones de negocio: balance y transfer}

\subsubsection*{Consulta de saldo}

\codigo{api/BalanceEndpoint.java} atiende \codigo{GET /api/accounts/\{id\}} y
devuelve la instantánea de una cuenta. Exige un JWT válido: lo primero que hace,
tras descartar métodos distintos de \codigo{GET} (\codigo{405}), es llamar a
\codigo{auth.authorize(req)}; si regresa \codigo{null} responde \codigo{401}.
Luego resuelve la variable de ruta \codigo{id} con \codigo{req.pathParam("id")}:
si no es un entero (no encaja en el patrón \codigo{-?\textbackslash{}d+})
responde \codigo{400} con \codigo{\{"error":"bad\_id"\}}; si es numérico pero se
sale del rango de \codigo{int}, o si la cuenta no está registrada, responde
\codigo{404}. Con una cuenta válida construye el cuerpo con la forma obligatoria
del contrato, donde las llaves son exactamente \codigo{id}, \codigo{propietario}
y \codigo{balance}, por ejemplo
\codigo{\{"id":125,"propietario":"JUAN MOLINAR HERNANDEZ","balance":15750.25\}}.
Un detalle clave: \codigo{id} y \codigo{balance} deben serializarse como
\emph{números} JSON (sin comillas); como \codigo{Money.toDecimal()} devuelve un
\codigo{String}, el saldo se envuelve en un \codigo{BigDecimal} para que
\codigo{Gson} lo emita como número y no como cadena entre comillas.

\subsubsection*{Transferencia}

\codigo{api/TransferEndpoint.java} atiende \codigo{POST /api/transactions/transfer}
y mueve dinero entre dos cuentas. También exige JWT (\codigo{401} sin
\emph{token}) y rechaza métodos distintos de \codigo{POST} (\codigo{405}). El
cuerpo obligatorio usa la forma
\codigo{\{"sourceAccountId":"123","targetAccountId":"456","amount":200.00\}},
donde los identificadores son \emph{cadenas} y \codigo{amount} es decimal. El
monto decimal se convierte a centavos con \codigo{Money.toCents()} y la operación
se orquesta en \codigo{Bank.transfer()}, que valida la solicitud, aplica el
movimiento atómico, estampa un número de secuencia fresco y preserva la
invariante de conservación. Un éxito devuelve \codigo{200} con
\codigo{\{"status":"ok","seq":<seq>\}} llevando la secuencia asignada. Un rechazo
se manifiesta como una \codigo{TransferException}: su \codigo{code()} determina el
estado, donde \codigo{"no\_such\_account"} mapea a \codigo{404} y el resto de
códigos (\codigo{"self\_transfer"}, \codigo{"bad\_amount"}, \codigo{"low\_balance"})
mapean a \codigo{400}, y el cuerpo refleja \codigo{\{"error":<code>\}}. Un JSON
mal formado, un identificador no numérico o un monto inválido producen
\codigo{400} con \codigo{\{"error":"bad\_request"\}}.

\begin{lstlisting}[language=none,caption={Contrato JSON de POST /api/transactions/transfer}]
POST /api/transactions/transfer
Authorization: Bearer <jwt>
Content-Type: application/json

{"sourceAccountId":"123","targetAccountId":"456","amount":200.00}

200 OK
{"status":"ok","seq":42}

400 Bad Request
{"error":"low_balance"}

404 Not Found
{"error":"no_such_account"}
\end{lstlisting}

El mapeo de \codigo{TransferException} a estado HTTP se concentra en pocas
líneas, copiadas tal cual del endpoint:

\begin{lstlisting}[language=Java,caption={Mapeo de rechazos de transferencia a estado REST en TransferEndpoint.java}]
} catch (TransferException e) {
    int status = "no_such_account".equals(e.code()) ? 404 : 400;
    JsonObject err = new JsonObject();
    err.addProperty("error", e.code());
    return Reply.json(status, gson.toJson(err));
}
\end{lstlisting}

\subsection{Endpoints de monitoreo}

Tres endpoints alimentan el tablero. Aquí se resumen; su funcionamiento detallado
se trata en el capítulo 9.

\codigo{api/StatsEndpoint.java} atiende \codigo{GET /api/stats} y reporta las
métricas de \emph{este} nodo. Serializa la instantánea viva de
\codigo{NodeStats} llamando a \codigo{stats.toJson()}, que produce un objeto JSON
con las llaves \codigo{nodeId}, \codigo{role}, \codigo{status},
\codigo{accountCount}, \codigo{totalBalance}, \codigo{transferCount},
\codigo{lastTxId}, \codigo{cpuPercent}, \codigo{ramPercent}, \codigo{diskPercent}
y \codigo{gcsCount} (el número de transacciones almacenadas en Cloud Storage).
Responde \codigo{405} ante un método distinto de \codigo{GET}.

\codigo{api/PanelEndpoint.java} atiende \codigo{GET /panel} y agrega las métricas
de todo el clúster. Combina la instantánea local (\codigo{stats.toJson()}) con las
métricas que obtiene por HTTP del endpoint \codigo{/api/stats} de cada par
listado en \codigo{TES\_PEERS} (vía \codigo{NodeConfig}), y produce un único
documento JSON con la forma \codigo{\{"nodes":[...]\}}, una entrada por nodo, para
que el tablero dibuje el clúster en una sola página. Cada par se consulta de forma
concurrente con un \codigo{HttpClient} de \emph{timeout} acotado; un par
inalcanzable o lento no tumba la respuesta: se reporta como nodo caído
(\codigo{"status":"down"}) en lugar de fallar todo el panel. Además recuerda el
último \codigo{nodeId} que cada par reportó mientras estuvo arriba, de modo que la
tarjeta de un nodo caído conserve su identidad real durante la demostración de
tolerancia a fallos.

\codigo{api/DashboardEndpoint.java} atiende \codigo{GET /} y sirve el tablero de
una sola página: transmite el recurso \codigo{dashboard.html} desde el
\emph{classpath} como respuesta HTML. La propia página sondea \codigo{/panel} (y
ofrece un control de \emph{Refresh}) para mostrar el estado de cada nodo, su
número de cuentas, el saldo total, el número de transferencias, el id de la
última transacción, \%CPU, \%RAM, \%Disco y el conteo de transacciones en Cloud
Storage. Responde \codigo{405} ante un método distinto de \codigo{GET} y
\codigo{404} si el recurso HTML no existe.

\section{Replicacion y reanudacion exacta por secuencia}

Este capitulo es el punto clave de la rubrica de tolerancia a fallos: cuando se
mata el \emph{proceso} de una replica (no la maquina virtual) y luego revive, la
replica debe reanudar \textbf{desde el numero de secuencia exacto} en el que se
quedo, pidiendo al lider solo el delta faltante. El anti-patron de ``borrar el
estado y volver a descargar todo'' funciona pero no escala y cuesta puntos. Las
piezas que lo implementan son \codigo{cluster/ReplicaFeed.java} (lado lider),
\codigo{cluster/ReplicaSync.java} (lado replica), \codigo{cluster/WireCodec.java}
(formato de mensaje) y \codigo{durable/ReplicaCheckpoint.java} (durabilidad que
sobrevive a la muerte del proceso).

\subsection{Protocolo de replicacion}

El lider publica su flujo de \emph{commits} sobre TCP plano a traves de
\codigo{cluster/ReplicaFeed.java}, que escucha en el puerto de replicacion
\codigo{TES\_REPL\_PORT}. Cuando una replica se conecta, primero anuncia hasta
donde ya esta al dia con un saludo de una linea \codigo{CATCHUP <seq>}; el lider
lo lee en \codigo{serveReplica()} y \codigo{parseCatchup()}. La respuesta consta
de dos fases: primero reproduce, en orden, todas las transferencias del
\codigo{core/TransferLog.java} cuya secuencia es estrictamente mayor que esa marca
(\codigo{log.since(since)}), y despues sigue transmitiendo en vivo cada nuevo
commit.

El detalle decisivo es que \textbf{el lider no espera confirmacion de las
replicas}: nunca bloquea su throughput por una replica lenta o caida.
\codigo{ReplicaFeed} es un \codigo{core/CommitListener.java} invocado en linea
sobre los hilos de commit del lider, asi que no puede bloquearse ahi. Cada
suscriptor tiene su propia cola acotada (\codigo{QUEUE\_CAPACITY = 8192}) y un
hilo escritor dedicado; \codigo{onCommit()} solo hace un \codigo{offer} no
bloqueante en cada cola. Una replica cuyo socket se atasca solo llena su
\emph{propia} cola; al desbordarse, se descarta ese suscriptor y se cierra su
socket, sin afectar la ruta de commit ni a las demas replicas.

\begin{lstlisting}[language=Java,caption={ReplicaFeed.onCommit: fan-out no bloqueante hacia las replicas}]
@Override
public void onCommit(Transfer t) {
    String line = WireCodec.encode(t);
    for (Subscriber subscriber : subscribers) {
        if (!subscriber.enqueue(line)) {
            drop(subscriber);
        }
    }
}
\end{lstlisting}

El formato de cada mensaje lo define \codigo{cluster/WireCodec.java}: una
transferencia se serializa como un unico objeto JSON por linea con los cuatro
campos que un par necesita para reaplicar el commit (la secuencia, las cuentas
origen y destino, y el monto en centavos). Al ser texto y un objeto por linea, la
replica puede leer el flujo con un lector de lineas ordinario. \codigo{WireCodec}
es un ayudante sin estado (constructor privado) que solo expone
\codigo{encode()} y \codigo{decode()}; los campos numericos se leen con su ancho
nativo (\codigo{getAsLong}/\codigo{getAsInt}) para que \codigo{seq} y
\codigo{cents} de 64 bits sobrevivan el viaje de ida y vuelta sin perder
precision.

\begin{lstlisting}[language=Java,caption={WireCodec: codificar y decodificar una transferencia como linea JSON}]
public static String encode(Transfer t) {
    JsonObject object = new JsonObject();
    object.addProperty("seq", t.seq());
    object.addProperty("from", t.from());
    object.addProperty("to", t.to());
    object.addProperty("cents", t.cents());
    return object.toString();
}

public static Transfer decode(String line) {
    try {
        JsonObject object = JsonParser.parseString(line).getAsJsonObject();
        long seq = object.get("seq").getAsLong();
        int from = object.get("from").getAsInt();
        int to = object.get("to").getAsInt();
        long cents = object.get("cents").getAsLong();
        return new Transfer(seq, from, to, cents);
    } catch (RuntimeException ex) {
        throw new IllegalArgumentException("malformed transfer line: " + line, ex);
    }
}
\end{lstlisting}

La linea de cable es entonces \codigo{\{"seq":..,"from":..,"to":..,"cents":..\}}
seguida de un salto de linea. El orden se mantiene sin un candado global de
difusion: el lider agrega la transferencia al log \emph{antes} de notificar a los
oyentes, y un suscriptor se registra en la audiencia en vivo \emph{antes} de que
su escritor lea el backlog, de modo que todo commit llega a la replica por el
backlog, por la cola en vivo, o (inofensivamente) por ambos. Los duplicados y el
ligero reordenamiento en esa frontera los absorbe la replica.

\subsection{La replica: aplicar y seguir el flujo}

Del lado seguidor, \codigo{cluster/ReplicaSync.java} mantiene al nodo alineado con
el lider. Marca al puerto de replicacion del \codigo{TES\_LEADER\_HOST} y abre la
conversacion con el saludo \codigo{CATCHUP <seq>}, donde la secuencia es
\codigo{Bank.lastSeq()}, la transferencia mas alta que el \codigo{core/Bank.java}
local ya aplico. El lider entonces transmite, como lineas JSON, primero el backlog
posterior a esa marca y luego cada commit en vivo; cada linea recibida se vuelve a
convertir en un \codigo{core/Transfer.java} con \codigo{WireCodec.decode()}.

Como el lider reparte los commits desde hilos concurrentes, una secuencia mas alta
puede llegar antes que una mas baja. Por eso \codigo{ReplicaSync} aplica las
transferencias a traves de un pequeno \emph{buffer de reordenamiento}
(\codigo{TreeMap<Long, Transfer> pending}) que garantiza que el banco siempre las
vea en orden de secuencia estricto y sin huecos. \codigo{bufferAndDrain()}
guarda las llegadas fuera de orden hasta que las secuencias faltantes se
completan, y entonces las vacia de forma contigua hacia
\codigo{Bank.applyReplicated()}. Una transferencia con secuencia igual o menor a
la marca del banco es un duplicado y se descarta, asi que un solapamiento
reenviado tras una reconexion es inofensivo.

\begin{lstlisting}[language=Java,caption={ReplicaSync.bufferAndDrain: reordenar y aplicar en secuencia contigua}]
private void bufferAndDrain(TreeMap<Long, Transfer> pending, Transfer t) {
    if (t.seq() > bank.lastSeq()) {
        pending.put(t.seq(), t);
    }
    while (!pending.isEmpty() && pending.firstKey() <= bank.lastSeq() + 1) {
        Transfer next = pending.pollFirstEntry().getValue();
        if (next.seq() > bank.lastSeq()) {
            bank.applyReplicated(next);
        }
    }
}
\end{lstlisting}

El enlace se trata como poco confiable. El lector corre dentro de
\codigo{syncLoop()}, un bucle que, ante cualquier caida, espera brevemente
(\codigo{RECONNECT\_DELAY\_MS = 1000}) y vuelve a marcar, reemitiendo
\codigo{CATCHUP} con la marca mas reciente del banco para que la reproduccion se
reanude desde ahi y no desde cero. Cada reconexion estrena un buffer
\codigo{pending} nuevo, de modo que cualquier resto fuera de orden del enlace
anterior se descarta. Una linea corrupta tampoco mata al lector: rompe el ciclo
interno para que el bucle externo reconecte y el lider reenvie desde la marca
actual.

\subsection{Checkpoint local y reanudacion exacta}

Todo lo anterior basta para sobrevivir a una caida de \emph{red}: la marca vive en
memoria en \codigo{Bank.lastSeq()}. Pero si matan el \emph{proceso} de la replica,
esa memoria se pierde. Aqui entra \codigo{durable/ReplicaCheckpoint.java}, que
persiste en almacenamiento durable lo necesario para reanudar exacto: la marca de
secuencia y los saldos. Un checkpoint es un unico objeto de Cloud Storage,
sobrescrito en el mismo lugar, con la forma
\codigo{\{v, watermark, idBase, count, balances\}}. El \codigo{watermark} es la
secuencia mas alta aplicada; \codigo{balances} es una instantanea posicional donde
el casillero \codigo{k} corresponde al saldo de la cuenta \codigo{idBase + k}, asi
que la posicion codifica el id y no se guarda una clave por cuenta. La marca y los
saldos se capturan juntos bajo el monitor del banco (\codigo{Bank.snapshotInto()}),
de modo que el par es consistente aunque el hilo de sincronizacion este aplicando
transferencias.

Al revivir, \codigo{load()} se ejecuta \emph{antes} de que arranque la
replicacion: restaura los saldos sobre el ledger ya cargado del CSV y siembra la
marca del banco con \codigo{Bank.seedLastSeq(watermark)}, que sube (nunca baja) el
\codigo{lastSeq} y el contador de secuencia. Asi, el \codigo{CATCHUP} posterior
lleva la marca restaurada \codigo{X} en vez de \codigo{0}, y el lider solo
transmite el delta, es decir desde \codigo{X+1} (\codigo{log.since(X)} devuelve las
transferencias \emph{estrictamente mayores} que \codigo{X}). Un checkpoint
ausente, ilegible o incompatible (version, \codigo{idBase} o numero de cuentas que
no coinciden) se ignora y el arranque cae a un catch-up completo; \codigo{load()}
nunca lanza excepcion.

El mensaje que evidencia la reanudacion exacta -- el literal ``Me quede en la
secuencia X, enviame desde X+1'' que pide la rubrica -- lo imprime
\codigo{cluster/ReplicaSync.java} justo antes de cada \codigo{CATCHUP}, en el camino
de revival de toda replica (con o sin checkpoint). Captura la marca una sola vez
para que el mensaje y el \codigo{CATCHUP} coincidan:

\begin{lstlisting}[language=Java,caption={ReplicaSync: imprime el punto de reanudacion antes del CATCHUP}]
long seq = bank.lastSeq();
System.out.println("[replica] Me quede en la secuencia " + seq
        + ", enviame desde " + (seq + 1));
out.write(("CATCHUP " + seq + "\n").getBytes(StandardCharsets.UTF_8));
\end{lstlisting}

Cuando ademas hay checkpoint, \codigo{ReplicaCheckpoint.load()} deja constancia de
la marca restaurada y \codigo{app/Node.java} la confirma, de modo que la salida
tipica al revivir es:

\begin{lstlisting}[language=none,caption={Salida tipica al revivir el proceso de una replica}]
[checkpoint] restored checkpoint/nodo-2.json at watermark 4127
[node] replica resuming at sequence 4127
[replica] Me quede en la secuencia 4127, enviame desde 4128
\end{lstlisting}

Con la marca \codigo{4127} sembrada, \codigo{ReplicaSync} envia
\codigo{CATCHUP 4127} y el lider reproduce desde la \codigo{4128} en adelante: esa
es la reanudacion exacta por secuencia. La cadencia de guardado tiene dos partes:
un paro ordenado (el \emph{shutdown hook} en \codigo{Node.java}) escribe el estado
exacto de la ultima transferencia aplicada -- y la demostracion en vivo apaga la VM,
que es un paro ordenado --, mientras que un demonio periodico
(\codigo{ReplicaCheckpoint.start()}) escribe un piso para que un \emph{crash} duro
pierda a lo sumo un intervalo. Un \codigo{save()} cuya marca no cambio desde el
ultimo se omite (\codigo{lastSavedWatermark}), asi que una replica ociosa no cuesta
nada.

\subsection{Por que no ``borra y descarga''}

El contraste con el anti-patron es el corazon de la rubrica. ``Borrar el estado y
bajar todo de nuevo'' significa que la replica, al revivir, descarta sus saldos y
pide al lider la historia completa desde la secuencia \codigo{0}. Eso funciona,
pero el costo de catch-up crece con \emph{toda} la historia del banco, no con lo
que la replica perdio mientras estuvo caida: no escala y, segun la rubrica, baja la
calificacion de 25 a 13 puntos.

La implementacion aqui reanuda por \emph{delta exacto}. Gracias al checkpoint
durable, la replica conoce su \codigo{watermark} al arrancar y solo pide
\codigo{CATCHUP <watermark>}; el lider responde con \codigo{log.since(watermark)},
que son unicamente las transferencias \codigo{watermark+1, watermark+2, ...} que
ocurrieron durante la caida. El trabajo de recuperacion es proporcional a la
\emph{brecha}, no al tamano total del estado. Y como \codigo{applyReplicated()}
descarta cualquier secuencia ya aplicada y el buffer de reordenamiento garantiza
orden contiguo, reanudar desde la marca es seguro aunque el lider reenvie un
solapamiento: el resultado converge al estado del lider sin reprocesar la historia
entera.

\section{Persistencia durable en Cloud Storage}

La durabilidad del mini-banco se apoya en un \emph{log} de transferencias que
vive en un bucket de Google Cloud Storage. Cada transferencia confirmada por el
lider termina escrita en esa bitacora, de modo que si caen los tres nodos a la
vez el lider pueda renacer y reconstruir todos los saldos releyendo el log. Tres
clases cooperan: \codigo{durable/Journal.java} orquesta la escritura y la
recuperacion, \codigo{durable/GcsStore.java} habla con la JSON API de Cloud
Storage y \codigo{durable/GcsAuth.java} obtiene los tokens de acceso. La unidad
que se persiste es siempre un \codigo{core/Transfer.java}, el mismo registro que
viaja por el cluster.

\subsection{La bitacora de transacciones}

El \codigo{durable/Journal.java} es un orquestador de durabilidad que solo existe
en el lider (las replicas nunca poseen un \codigo{Journal}). Implementa la
interfaz \codigo{core/CommitListener.java}: una vez registrado sobre el banco del
lider, su metodo \codigo{onCommit(Transfer)} recibe cada transferencia recien
confirmada. Lo importante es que \emph{no} sube nada a GCS en ese momento: solo
encola la transferencia y regresa de inmediato. Asi, la respuesta HTTP de una
transferencia nunca se bloquea esperando un viaje de ida y vuelta a Cloud
Storage.

\begin{lstlisting}[language=Java,caption={onCommit solo encola; la subida ocurre fuera del hot path}]
@Override
public void onCommit(Transfer transfer) {
    queue.offer(transfer);
}
\end{lstlisting}

La cola es un \codigo{BlockingQueue} que un unico hilo de fondo
(\codigo{journal-writer}, marcado como \emph{daemon}) consume en
\codigo{drainLoop()}. Ese hilo espera con \codigo{poll(POLL\_MS, ...)} a que
aparezca una transferencia, y cuando la toma vacia de golpe todo lo que haya
disponible en la cola hasta un maximo de \codigo{MAX\_BATCH} (1000) con
\codigo{drainTo}, componiendo un solo lote. El lote completo se sube como un
unico objeto mediante \codigo{store.putBatch(batch)} y, solo si la subida tuvo
exito, se avanza el contador \codigo{stored} en el tamano del lote. El agrupado
por lotes es lo que permite que el log durable siga el ritmo de una tasa alta de
commits sin un objeto por transferencia.

\begin{lstlisting}[language=Java,caption={drainLoop compone un lote y lo sube como un solo objeto}]
batch.clear();
batch.add(first);
queue.drainTo(batch, MAX_BATCH - 1);
try {
    store.putBatch(batch);
    stored.addAndGet(batch.size());
} catch (IOException e) {
    System.err.println("journal: could not persist batch of " + batch.size()
            + " ending at seq " + batch.get(batch.size() - 1).seq()
            + ": " + e.getMessage());
}
\end{lstlisting}

Cada entrada del log es exactamente un \codigo{core/Transfer.java}: un registro
inmutable con un numero de secuencia monotono asignado por el ledger al confirmar
(\codigo{seq}), la cuenta origen (\codigo{from}), la cuenta destino
(\codigo{to}) y el monto movido en centavos (\codigo{cents}). Es la misma forma
que serializan el codec del cluster y el journal de GCS.

\begin{lstlisting}[language=Java,caption={La unidad persistida: el registro Transfer (core/Transfer.java)}]
public record Transfer(long seq, int from, int to, long cents) {
}
\end{lstlisting}

El contador \codigo{stored()} reporta cuantas transferencias estan
\emph{realmente} subidas a GCS; queda por debajo del numero de commits aplicados
en exactamente la profundidad de la cola, que es la cifra durable que debe
mostrar el dashboard. Un fallo de subida se registra en \codigo{stderr} pero no
detiene el hilo: el lazo corre mientras \codigo{running} sea verdadero o quede
algo en la cola. Un lote cuya subida falla se descarta y deja un hueco en el log
durable; para que ese hueco no pase inadvertido, \codigo{recover()} valida la
contiguidad al arrancar (el log es contiguo desde 1) y avisa por \codigo{stderr} si
una secuencia no es la sucesora esperada, en vez de recuperar en silencio un log
incompleto.

Esta estrategia asincrona tiene una ventana de durabilidad explicita. En un
apagado ordenado, \codigo{stop()} pone \codigo{running} en falso y hace
\codigo{join} al escritor para drenar la cola, de modo que toda transferencia ya
confirmada quede subida antes de que el proceso termine; \codigo{app/Node.java}
registra ese \codigo{stop()} como \emph{shutdown hook}. En cambio, una muerte
abrupta del proceso deja la cola en vuelo como la ventana de durabilidad inherente
a un journaling asincrono: como el commit responde al cliente antes de subir a
GCS, una transferencia confirmada puede perderse en un \emph{kill} duro. Ese es
el precio consciente de mantener Cloud Storage fuera del camino critico.

\subsection{Acceso a Cloud Storage}

El \codigo{durable/GcsStore.java} es un cliente delgado sobre la JSON API REST de
Cloud Storage, construido directamente con \codigo{java.net.http.HttpClient} y
sin el SDK de Google. Todos los objetos del journal cuelgan del prefijo
\codigo{journal/}. El metodo \codigo{put(Transfer)} existe y escribiria una
transferencia suelta como \codigo{journal/tx-seq.json}, pero el camino vivo no lo
usa: \codigo{Journal} siempre llama a \codigo{putBatch}, que compone muchas
transferencias en un solo objeto \codigo{journal/batch-min-max.json} (un arreglo
JSON), nombrado por la secuencia minima y maxima del lote. Cloud Storage no tiene
\emph{append}; componer objetos es la alternativa sancionada por el contrato.

Las subidas usan el endpoint de \emph{simple media upload}
(\codigo{UPLOAD\_BASE}) con \codigo{uploadType=media}, mientras que listar y leer
usan la JSON API (\codigo{API\_BASE}). Toda peticion lleva la cabecera
\codigo{Authorization: Bearer} con un token de \codigo{GcsAuth} y se serializa con
gson. Ademas, \codigo{GcsStore} ofrece \codigo{putCheckpoint}/\codigo{getCheckpoint}
para que una replica guarde su snapshot bajo el prefijo \codigo{checkpoint/};
estos quedan deliberadamente fuera de \codigo{journal/} para que el
\codigo{readAll()} del lider nunca los liste ni intente parsearlos como
\codigo{Transfer}.

\begin{lstlisting}[language=Java,caption={putBatch compone el lote en un solo objeto journal/batch-min-max.json}]
long min = batch.stream().mapToLong(Transfer::seq).min().getAsLong();
long max = batch.stream().mapToLong(Transfer::seq).max().getAsLong();
String name = PREFIX + "batch-" + min + "-" + max + ".json";
String url = UPLOAD_BASE + bucket + "/o?uploadType=media&name=" + enc(name);
HttpRequest request = HttpRequest.newBuilder(URI.create(url))
        .timeout(HTTP_TIMEOUT)
        .header("Authorization", "Bearer " + auth.accessToken())
        .header("Content-Type", "application/json")
        .POST(HttpRequest.BodyPublishers.ofString(gson.toJson(batch)))
        .build();
\end{lstlisting}

Toda peticion pasa por \codigo{sendWithRetry}, que reintenta hasta
\codigo{MAX\_ATTEMPTS} (3) veces ante un fallo de red transitorio
(\codigo{IOException}: una conexion caida, un timeout) y, si todos los intentos
fallan, relanza la excepcion. Esto hace que la recuperacion en frio aborte de
forma ruidosa en lugar de servir silenciosamente un estado solo-base. Un estado
HTTP no-2xx no se reintenta: se devuelve al llamador, que decide como tratarlo
(por ejemplo, \codigo{getCheckpoint} interpreta un 404 como ``aun no existe'' y
devuelve \codigo{null}).

\subsubsection{Autenticacion con la cuenta de servicio}

El \codigo{durable/GcsAuth.java} acuna tokens de acceso OAuth2 para la JSON API
usando el flujo de \emph{key-file} de cuenta de servicio, tambien sin SDK de
Google (el metadata server de la VM no se usa a proposito). En el constructor lee
el JSON de la cuenta de servicio apuntado por \codigo{TES\_GCS\_KEYFILE} y extrae
\codigo{client\_email}, \codigo{token\_uri} (con un endpoint por defecto si falta)
y \codigo{private\_key}, del que reconstruye la llave RSA PKCS8.

Para obtener un token, \codigo{buildAssertion()} arma un JWT con cabecera RS256 y
un claim set cargado con \codigo{iss}, \codigo{scope}
(\codigo{devstorage.read\_write}), \codigo{aud}, \codigo{iat} y \codigo{exp}, lo
firma con \codigo{SHA256withRSA} sobre la llave privada y lo codifica en Base64
URL sin relleno. No se envia \codigo{sub} (eso es solo para delegacion de
dominio y aqui seria rechazado como \codigo{unauthorized\_client}).

\begin{lstlisting}[language=Java,caption={Claims del JWT que GcsAuth autofirma (RS256)}]
JsonObject claims = new JsonObject();
claims.addProperty("iss", clientEmail);
claims.addProperty("scope", SCOPE);
claims.addProperty("aud", tokenUri);
claims.addProperty("iat", now);
claims.addProperty("exp", now + TOKEN_TTL_SECONDS);
\end{lstlisting}

El JWT firmado se intercambia en \codigo{refresh()} por un POST
\codigo{x-www-form-urlencoded} al \codigo{token\_uri} con el grant
\codigo{urn:ietf:params:oauth:grant-type:jwt-bearer}, y del cuerpo se lee
\codigo{access\_token}. \codigo{accessToken()} cachea el token y lo reusa hasta
poco antes de expirar (margen \codigo{REFRESH\_SKEW\_SECONDS} de 60 s), acunando
uno nuevo solo cuando hace falta; el metodo es \codigo{synchronized} porque el
hilo escritor lo consulta en cada subida.

El alta de la cuenta de servicio se hace una sola vez, segun describe
\codigo{SETUP-GCS.md}: crear la cuenta, otorgarle \codigo{roles/storage.objectAdmin}
sobre el bucket (no basta \codigo{objectCreator}, porque el log no solo sube sino
que \emph{relee} en el arranque en frio) y descargar la llave
\codigo{credentials.json}. Esa llave nunca se sube al repositorio; en la VM se
copia aparte y se apunta \codigo{TES\_GCS\_KEYFILE} a su ruta, mientras
\codigo{TES\_BUCKET} lleva el nombre del bucket (solo en el nodo lider).

\subsection{Recuperacion ante desastre total}

Cuando el lider arranca, \codigo{app/Node.java} corre la recuperacion durable
antes de aceptar trafico, siempre que Cloud Storage este configurado (si faltan
\codigo{TES\_BUCKET} o \codigo{TES\_GCS\_KEYFILE}, el lider sigue sirviendo todos
los endpoints en memoria y el journal simplemente no se activa). El arranque crea
el \codigo{GcsAuth} y el \codigo{GcsStore}, instancia el \codigo{Journal} y llama
a \codigo{recover}, pasandole \codigo{bank::applyReplicated} como la accion que
aplica cada transferencia recuperada al ledger.

\begin{lstlisting}[language=Java,caption={Arranque del lider: recuperar, sembrar el contador y registrar el journal (app/Node.java)}]
GcsAuth auth = new GcsAuth(Path.of(config.gcsKeyfile()));
GcsStore store = new GcsStore(auth, config.bucket());
journal = new Journal(store);
int recovered = journal.recover(bank::applyReplicated);
journal.start(recovered);
bank.addCommitListener(journal);
\end{lstlisting}

Dentro de \codigo{recover(CommitListener apply)}, el journal pide a
\codigo{store.readAll()} todas las transferencias persistidas y reaplica cada una
en orden al ledger. El metodo \codigo{readAll()} de \codigo{GcsStore} lista los
objetos bajo el prefijo \codigo{journal/} (siguiendo la paginacion por
\codigo{nextPageToken} de la JSON API), baja cada uno y los aplana: si el objeto
es un arreglo (un lote) expande cada elemento, y si es un objeto suelto lo
convierte directo. Como cada \codigo{Transfer} carga su \codigo{seq}, el ``donde
quedo'' no depende del nombre de los objetos ni de un cursor persistido del
lider: \codigo{readAll()} reordena todo por secuencia al final, asi que el ledger
se reconstruye exactamente hasta la ultima secuencia confirmada sin huecos ni
duplicados de orden.

\begin{lstlisting}[language=Java,caption={readAll aplana lotes y transferencias sueltas, y reordena por secuencia (GcsStore)}]
JsonElement parsed = JsonParser.parseString(response.body());
if (parsed.isJsonArray()) {
    for (JsonElement element : parsed.getAsJsonArray()) {
        transfers.add(gson.fromJson(element, Transfer.class));
    }
} else {
    transfers.add(gson.fromJson(parsed, Transfer.class));
}
// ...
transfers.sort(Comparator.comparingLong(Transfer::seq));
\end{lstlisting}

Hecha la reaplicacion, \codigo{recover} devuelve cuantas transferencias se
recuperaron y \codigo{start(recovered)} siembra con ese valor el contador
\codigo{stored} (para que el dashboard arranque mostrando lo ya durable) y lanza
el hilo escritor que sube los commits vivos a partir de ese punto. El lider queda
entonces sirviendo con los saldos completos reconstruidos desde el log, que el
dashboard muestra como saldos recuperados. La garantia es fuerte para todo lo que
alcanzo a subirse a GCS y se vio reforzada por el \emph{shutdown hook} que drena
la cola en un apagado ordenado; lo unico no recuperable es la ventana asincrona
en vuelo descrita antes, ante una caida abrupta.

\section{Dashboard de observabilidad}

El dashboard es la cara visible del cluster: una sola página web que monitorea
los tres nodos a la vez y, para cada uno, pinta los datos de negocio (cuentas,
saldo, transferencias, secuencia de la última transacción) junto con las
métricas de infraestructura del sistema operativo de su máquina (CPU, RAM y
Disco). El armado descansa en tres endpoints HTTP, registrados en
\codigo{app/Node.java} sobre el router \codigo{net/Routes.java}: \codigo{GET /}
sirve la página, \codigo{GET /panel} agrega la vista de todo el cluster y
\codigo{GET /api/stats} expone las métricas de un solo nodo. El extra de la
rúbrica (auto-refresh, sin tener que pulsar ``Refrescar'') se cubre con un
\codigo{setInterval} de polling en el JavaScript de la página.

\subsection{La interfaz web}

La página vive como recurso estático en \codigo{src/main/resources/dashboard.html}
y se entrega tal cual desde el classpath por \codigo{api/DashboardEndpoint.java}:
ese endpoint solo valida que el método sea \codigo{GET} (responde \codigo{405} si
no) y vuelca el HTML con \codigo{Reply.html(200, ...)}, leyendo el recurso
\codigo{/dashboard.html} en cada petición. No tiene colaboradores ni lógica de
negocio; toda la inteligencia está en el JavaScript del propio documento.

Es una SPA en HTML/CSS/JS puro, sin dependencias ni librerías externas. Al
cargar, y cada vez que se refresca, la función \codigo{loadPanel()} hace
\codigo{fetch("/panel")} pidiendo JSON y, con la respuesta, repinta tres zonas:
el \emph{hero} (saldo total del cluster, número de cuentas y última transacción,
tomados del nodo líder), un sello de consistencia que compara los saldos totales
de los nodos en línea para indicar si el invariante se conserva, y una tarjeta
por nodo. Cada tarjeta dibuja el rol (\codigo{Lider}/\codigo{Replica}), el estado
(\codigo{Activo}/\codigo{Caido}), el saldo del nodo, los cuatro contadores de
negocio (cuentas, transferencias, última tx y transacciones en Cloud Storage) y
tres barras (\emph{gauges}) para CPU, RAM y Disco. Si \codigo{/panel} falla, se
muestra un \emph{banner} de error en vez de romper la página.

El auto-refresh, que da los 15 puntos extra, es un \emph{toggle} ``En vivo'' que
\textbf{arranca encendido por defecto}: el checkbox lleva \codigo{checked} y, al
cargar, la página llama \codigo{toggleReactive(checkbox.checked)}, de modo que el
tablero empieza a refrescarse solo \emph{sin pulsar ``Refrescar''}, que es justo lo
que pide el extra. \codigo{toggleReactive()} arranca un \codigo{setInterval} que
vuelve a llamar a \codigo{loadPanel()} cada \codigo{REACTIVE\_INTERVAL\_MS}
(1500 ms); al desactivar el toggle se limpia el temporizador con
\codigo{clearInterval}.

\begin{lstlisting}[language=none,caption={Polling automatico del dashboard (dashboard.html)}]
// Poll cadence (ms) used when the reactive toggle is enabled.
const REACTIVE_INTERVAL_MS = 1500;
let timer = null;
// ...
function toggleReactive(enabled) {
  if (timer) { clearInterval(timer); timer = null; }
  if (enabled) timer = setInterval(loadPanel, REACTIVE_INTERVAL_MS);
  loadPanel();
}
\end{lstlisting}

\subsection{Datos de negocio}

Los datos de negocio de cada nodo nacen en \codigo{app/NodeStats.java}, un
\emph{snapshot} en vivo que reparte las cifras entre dos colaboradores: del
\codigo{core/Ledger.java} salen el número de cuentas (\codigo{accountCount()} vía
\codigo{ledger.size()}) y el saldo total en centavos (\codigo{totalCents()}); del
\codigo{core/Bank.java} salen el número de transferencias confirmadas
(\codigo{transferCount()} vía \codigo{bank.appliedCount()}) y la secuencia de la
última transacción aplicada (\codigo{lastTxId()} vía \codigo{bank.lastSeq()}). El
método \codigo{toJson()} serializa todo a un \codigo{JsonObject} cuyas claves
coinciden exactamente con los campos que consume el JavaScript: \codigo{nodeId},
\codigo{role}, \codigo{status}, \codigo{accountCount}, \codigo{totalBalance}
(convertido a decimal con \codigo{Money.toDecimal()}), \codigo{transferCount},
\codigo{lastTxId}, los tres porcentajes y \codigo{gcsCount}.

Quien expone ese JSON de un solo nodo es \codigo{api/StatsEndpoint.java} en la
ruta \codigo{GET /api/stats}: valida el método y devuelve
\codigo{Reply.json(200, stats.toJson())}. Sobre él se monta
\codigo{api/PanelEndpoint.java} (ruta \codigo{GET /panel}), que arma la vista del
cluster completo: parte del \codigo{toJson()} del nodo local y, en paralelo,
consulta el \codigo{/api/stats} de cada peer listado en \codigo{TES\_PEERS}
(\codigo{NodeConfig}). Las peticiones se lanzan concurrentes con
\codigo{CompletableFuture} para que un nodo caído sume a lo sumo un \emph{timeout}
a la latencia del panel, no uno por peer; un peer inalcanzable no rompe la
respuesta, sino que se reporta como nodo en estado \codigo{down}. El resultado es
un único documento \codigo{\{"nodes":[...]\}} con una entrada por nodo.

\begin{lstlisting}[language=Java,caption={PanelEndpoint agrega el nodo local y los peers}]
StringBuilder nodes = new StringBuilder();
nodes.append("{\"nodes\":[").append(stats.toJson());
for (CompletableFuture<String> peer : pending) {
    nodes.append(',').append(peer.join());
}
nodes.append("]}");
return Reply.json(200, nodes.toString());
\end{lstlisting}

\subsection{Métricas de infraestructura del SO}

Las tres barras de CPU, RAM y Disco no son inventadas: \codigo{NodeStats} las lee
del sistema operativo de cada VM con Java puro, sin librerías externas. El CPU se
toma de \codigo{/proc/stat}: \codigo{cpuPercent()} suma los \emph{jiffies} de la
línea agregada \codigo{cpu} y calcula la fracción ocupada como
$1 - \Delta\text{idle}/\Delta\text{total}$ entre dos lecturas consecutivas (entre
dos refrescos del dashboard), guardando el muestreo previo bajo
\codigo{cpuLock} para que las peticiones concurrentes no se pisen. La RAM viene
de \codigo{/proc/meminfo} (\codigo{ramPercent()} compara \codigo{MemTotal} contra
\codigo{MemAvailable}). El Disco se calcula con \codigo{java.io.File} sobre la
raíz \codigo{/}: \codigo{diskPercent()} usa \codigo{getTotalSpace()} y
\codigo{getUsableSpace()}. Las tres pasan por \codigo{clamp()} al rango
$[0, 100]$ y por \codigo{round1()} a un decimal antes de serializarse.

\begin{lstlisting}[language=Java,caption={Lectura de \%RAM y \%Disco del SO (NodeStats)}]
public double ramPercent() {
    try {
        long total = readMeminfo("MemTotal");
        long available = readMeminfo("MemAvailable");
        if (total <= 0L) {
            return 0.0;
        }
        return clamp((double) (total - available) / total * 100.0);
    } catch (RuntimeException | IOException e) {
        return 0.0;
    }
}
\end{lstlisting}

Estos porcentajes viajan al tablero por la misma vía que los datos de negocio:
\codigo{toJson()} los empaqueta como \codigo{cpuPercent}, \codigo{ramPercent} y
\codigo{diskPercent}; \codigo{StatsEndpoint} los publica en \codigo{/api/stats},
\codigo{PanelEndpoint} los agrega en \codigo{/panel}, y el JavaScript de
\codigo{dashboard.html} los dibuja como \emph{gauges} (la barra se pinta en tono
de alarma cuando el valor llega al 80\%). Así, cada métrica del cuadro proviene
directamente del SO de la máquina que la reporta.

\section{Generador de carga, pruebas y despliegue}

Este capítulo cierra el manual con la cara operativa del sistema: cómo se le
inyecta carga para medirlo, cómo se prueba su lógica con un arnés en Java puro y
cómo se despliega y se opera el clúster sobre Google Cloud el día de la
evaluación.

\subsection{Generador de carga}

El generador de carga vive en \codigo{loadtest/LoadDriver.java}, una clase
\emph{stand-alone} que es el único ``cliente'' del proyecto: no depende de nada
del servidor más allá del contrato REST publicado. Genera tráfico contra un solo
nodo (el líder) durante una duración fija mezclando un \textbf{80\% de lecturas
de saldo y un 20\% de transferencias}, proporción gobernada por la constante
\codigo{READ\_RATIO = 0.80}. Por defecto corre \codigo{DEFAULT\_SECONDS = 60}
segundos con \codigo{DEFAULT\_CLIENTS = 8} hilos concurrentes, pero los
argumentos de \codigo{main} permiten fijar segundos, hilos y el rango de cuentas
\codigo{[minId, maxId]}; subiendo los segundos puede correr de forma
prácticamente indefinida, como exige la prueba del profesor que ``nunca
termina''.

Antes de medir, \codigo{bootstrapToken()} registra un usuario desechable
(\codigo{loaddriver}) y hace \codigo{login} para obtener el JWT que comparten
todos los hilos; tanto un \codigo{201} (creado) como un \codigo{409} (ya
registrado) son aceptables porque solo se necesita poder iniciar sesión. Un
detalle de rendimiento importante es que \emph{un} único \codigo{HttpClient}
(\codigo{sharedClient}) se reparte entre todos los \emph{workers}: un cliente por
hilo abriría cientos de hilos selectores propios y el generador se ahogaría en su
propio \emph{context-switch} antes de saturar al banco.

Cada hilo es un \codigo{Worker} que, hasta una fecha límite común en
\codigo{System.nanoTime()}, sortea la operación según \codigo{READ\_RATIO} y la
despacha. Solo se cuentan los éxitos reales: \codigo{doRead()} incrementa
\codigo{reads} únicamente si el servidor devolvió \codigo{HTTP 200} con un saldo
parseable, y \codigo{doTransfer()} incrementa \codigo{transfers} solo si la
transferencia fue aceptada (también \codigo{200}).

\begin{lstlisting}[language=Java,caption={Bucle de carga por hilo (Worker.run, recortado)}]
public void run() {
    HttpClient http = sharedClient;
    try {
        while (System.nanoTime() < deadlineNanos) {
            if (ThreadLocalRandom.current().nextDouble() < READ_RATIO) {
                doRead(http);
            } else {
                doTransfer(http);
            }
        }
    } finally {
        done.countDown();
    }
}
\end{lstlisting}

La parte didácticamente más valiosa es la verificación de la invariante de
conservación del dinero. Antes de arrancar la carga,
\codigo{captureTotalCents()} lee \codigo{GET /api/stats} una sola vez y guarda
\codigo{totalBalance} (en centavos): leer el total que el servidor suma en
memoria es O(1) en el cliente, en lugar de sumar cientos de miles de saldos por
HTTP. Al terminar, \codigo{verifyAndReport()} vuelve a leer el total y exige que
sea \textbf{idéntico} al inicial: ninguna transferencia crea ni destruye dinero,
así que la suma global debe permanecer constante. Si no coincide,
\codigo{reportCulprit()} recorre el rango buscando la cuenta con saldo negativo
para señalar al culpable; si coincide, imprime \codigo{CONSISTENT} junto con las
lecturas, las transferencias y el \codigo{score = transferencias*4 + lecturas}
(una transferencia vale cuatro veces una lectura, según el PDF del concurso).

\begin{lstlisting}[language=Java,caption={Verificacion de la invariante de conservacion (verifyAndReport, recortado)}]
void verifyAndReport(long initialTotalCents, String scenario) {
    long finalTotalCents = captureTotalCents();
    if (finalTotalCents != initialTotalCents) {
        System.out.println("INCONSISTENT: expected " + initialTotalCents
                + " cents, found " + finalTotalCents + " cents");
        reportCulprit(initialTotalCents, finalTotalCents);
        return;
    }
    long score = transfers.get() * 4 + reads.get();
    System.out.println("CONSISTENT");
    // ...
}
\end{lstlisting}

El cuerpo de la transferencia respeta la forma de cable obligatoria:
\codigo{transferBody()} serializa \codigo{sourceAccountId} y
\codigo{targetAccountId} como cadenas y \codigo{amount} como decimal; los montos
salen de \codigo{randomAmount()} en el rango \codigo{[0.01, 100.00]}, valores
bajos a propósito para que las transferencias rara vez agoten una cuenta y el
tráfico no degenere en puros rechazos por saldo insuficiente.

\subsection{Harness de pruebas}

El proyecto no usa JUnit: trae su propio arnés de pruebas en Java puro, formado
por dos piezas. \codigo{tests/Assert.java} es un ayudante mínimo de aserciones
(\codigo{isTrue}, \codigo{equals}, \codigo{notNull}, \codigo{throwsException})
donde cada método lanza un \codigo{AssertionError} con un mensaje descriptivo
cuando la expectativa falla. \codigo{tests/RunTests.java} es el punto de entrada:
instancia cada \emph{suite}, ejecuta sus casos y, capturando el
\codigo{AssertionError}, marca cada caso como \codigo{PASS} o \codigo{FAIL} sin
detener a los demás. Un caso es un \codigo{record Case(String name, Runnable
body)}, de modo que la falla se expresa simplemente lanzando una excepción desde
el cuerpo. Al final imprime un resumen \codigo{total/passed/failed} y termina con
estatus distinto de cero si algo falló, de manera que sirve para frenar la
compilación.

\begin{lstlisting}[language=Java,caption={Recoleccion y ejecucion de suites (RunTests.main, recortado)}]
List<Case> cases = new ArrayList<>();
cases.addAll(new LedgerTest().cases());
cases.addAll(new MoneyTest().cases());
cases.addAll(new RequestParserTest().cases());
cases.addAll(new WireCodecTest().cases());
cases.addAll(new AuthTest().cases());
cases.addAll(new ReplayTest().cases());
cases.addAll(new CheckpointTest().cases());
cases.addAll(new JwtCrossNodeTest().cases());
cases.addAll(new CatchupTest().cases());
// ... corre cada caso, cuenta PASS/FAIL y hace System.exit(1) si hubo fallos
\end{lstlisting}

Las \emph{suites} están organizadas por subsistema. \codigo{tests/MoneyTest.java}
verifica que el dinero se redondee y convierta correctamente entre cadena decimal
y centavos (con \codigo{HALF\_UP}) y que el viaje de ida y vuelta sea exacto.
\codigo{tests/LedgerTest.java} comprueba que \codigo{Ledger.move} transfiera
centavos exactos entre dos cuentas y mantenga el total constante bajo
concurrencia. \codigo{tests/RequestParserTest.java} valida que el parser arme una
petición tanto en una sola entrega como partida entre varias.
\codigo{tests/WireCodecTest.java} cubre la codificación de una transferencia a una
línea JSON y su \emph{round-trip}. \codigo{tests/AuthTest.java} prueba la
autorización del \codigo{Authenticator}: el esquema \codigo{Bearer} es
insensible a mayúsculas (RFC 7235) y tolera espacios, mientras que un encabezado
ausente, un esquema equivocado o un token vacío o falso se rechazan.

El bloque de mayor valor para este sistema distribuido son las pruebas de
recuperación y de réplica. \codigo{tests/ReplayTest.java} fija el contrato de
\codigo{Bank.applyReplicated()}: un seguidor reproduce en orden de secuencia y debe
\emph{forzar} la aplicación de cada transferencia (ya fue autorizada por el líder)
en vez de re-verificar fondos, porque re-verificar podría rechazar y
\emph{descartar} una transferencia autorizada mientras los contadores avanzan,
divergiendo de forma permanente; aplicar incondicionalmente es independiente del
orden y por tanto convergente. La \emph{suite} alimenta una secuencia en orden
adverso, comprueba la convergencia y que un reenvío de una secuencia ya vista sea
idempotente. \codigo{tests/CheckpointTest.java} cubre las
primitivas de \emph{checkpoint} que permiten a una réplica reanudar desde su
secuencia exacta tras un reinicio total: \codigo{Bank.snapshotInto} (captura
atómica de marca de agua y saldos), \codigo{Bank.seedLastSeq} (siembra la marca
desde un \emph{checkpoint} restaurado) y la garantía de que un \emph{snapshot}
restaurado más el delta del líder converge exactamente a los saldos por cuenta
del líder. \codigo{tests/CatchupTest.java} fija la reanudación por secuencia exacta:
que \codigo{log.since(X)} sea justamente la cola \codigo{X+1..} y que un banco
sembrado en la marca \codigo{X} descarte los reenvíos \codigo{<=X} y aplique desde
\codigo{X+1}. \codigo{tests/JwtCrossNodeTest.java} fija la propiedad entre nodos del
JWT: un token emitido con el secreto compartido valida en otro nodo y uno firmado
con un secreto distinto se rechaza.

Existen además dos verificaciones de integración que \textbf{no} corre
\codigo{RunTests} porque abren un socket real:
\codigo{tests/ReplConsistency.java} cablea un \codigo{Bank} líder con uno réplica
sobre \codigo{ReplicaFeed} y \codigo{ReplicaSync}, lo martilla con varios hilos
escritores haciendo miles de transferencias concurrentes y luego exige que el
saldo \textbf{por cuenta} líder $=$ réplica para cada cuenta (la suma global se
conserva aun si una transferencia se pierde o se duplica, así que solo la
comparación por cuenta atrapa un \emph{bug} de replicación). \codigo{tests/ReplInversion.java}
es la regresión de la convergencia por cuenta bajo carga concurrente con saldos
deliberadamente ajustados y montos de hasta el saldo completo -- la ruta que
originalmente expuso el \emph{bug} de inversión de secuencia --, que
\codigo{ReplConsistency} no ejercita porque sus cuentas arrancan demasiado holgadas. Cada una imprime \codigo{REPL\_CONSISTENCY\_OK} o
\codigo{REPL\_INVERSION\_OK} y sale con \codigo{0} en éxito; se ejecutan a mano
con el \emph{fat jar} en el \emph{classpath}.

\subsection{Despliegue en GCP}

La infraestructura se aprovisiona con \codigo{deploy/crear-infra.sh}, que crea en
Google Cloud (zona \codigo{us-central1-a}) una IP estática para el líder, las
reglas de \emph{firewall} (puerto \codigo{8080} abierto a clientes y
\emph{dashboard}, \codigo{9090} solo entre nodos con la etiqueta
\codigo{tesoreria}) y exactamente \textbf{2x \codigo{e2-standard-2} + 1x
\codigo{e2-standard-4}}, respetando la restricción de usar solo la serie E2. Los
tres nodos arrancan el \emph{mismo} \codigo{deploy/startup-script.sh} y difieren
únicamente por metadatos por nodo: \codigo{tes-node-id}, \codigo{tes-leader-host},
\codigo{tes-peers} y un \codigo{tes-jwt-secret} compartido. La convención clave es
que un \codigo{tes-leader-host} \emph{vacío} marca al nodo como líder (que recibe
la IP estática); las réplicas lo siguen por su nombre vía el DNS interno de la
VPC.

\begin{lstlisting}[language=none,caption={Creacion del cluster: 1 lider e2-standard-4 + 2 replicas e2-standard-2 (crear-infra.sh, recortado)}]
# Leader: empty tes-leader-host marks it the leader; it gets the static IP.
create_node "${LEADER}" "e2-standard-4" "" "${REPLICA_A}:${HTTP_PORT},${REPLICA_B}:${HTTP_PORT}" \
    --address="${LEADER_IP}"
# Replicas: follow the leader by name; e2-standard-2 each.
create_node "${REPLICA_A}" "e2-standard-2" "${LEADER}" "${LEADER}:${HTTP_PORT},${REPLICA_B}:${HTTP_PORT}"
create_node "${REPLICA_B}" "e2-standard-2" "${LEADER}" "${LEADER}:${HTTP_PORT},${REPLICA_A}:${HTTP_PORT}"
\end{lstlisting}

En cada arranque, \codigo{startup-script.sh} corre como \emph{root} y hace cuatro
cosas: instala el \emph{runtime} \codigo{openjdk-17-jre-headless}; descarga desde
Cloud Storage el \emph{fat jar}, el \emph{dataset} \codigo{alumnos.csv}, la unidad
\codigo{node.service} y la llave \codigo{gcs-key.json}; escribe el archivo de
entorno \codigo{/etc/tesoreria/node.env} con las variables \codigo{TES\_*}
leídas de los metadatos (entre ellas \codigo{TES\_NODE\_ID},
\codigo{TES\_PEERS}, \codigo{TES\_LEADER\_HOST}, \codigo{TES\_JWT\_SECRET} y
\codigo{TES\_BUCKET}); e instala y arranca la unidad \emph{systemd}. Que el mismo
script sirva para todos los roles es lo que mantiene la operación simple: la
diferencia líder/réplica se infiere de \codigo{TES\_LEADER\_HOST}.

La unidad \codigo{deploy/node.service} lanza el \emph{jar} con el puerto HTTP
como primer argumento (\codigo{ExecStart=/usr/bin/java -jar
.../tesoreria-distribuida.jar 8080}), lee el entorno de
\codigo{/etc/tesoreria/node.env} y reinicia \codigo{on-failure} para que un nodo
caído reingrese al clúster automáticamente. Lo \textbf{más importante} de esta
unidad es su política de apagado: usa \codigo{KillSignal=SIGTERM} con
\codigo{TimeoutStopSec=120} para que la JVM ejecute sus \emph{shutdown hooks}. Un
apagado \emph{graceful} con SIGTERM dispara, en una réplica, la escritura del
\emph{checkpoint} exacto (de ahí depende la garantía de reanudar desde la
secuencia precisa) y, en el líder, el drenado de la cola del \emph{journal}
durable. Por eso conviene distinguir reiniciar el \emph{proceso} de reiniciar la
\emph{VM}: matar el proceso con SIGTERM (o \codigo{systemctl stop}) preserva el
estado y el \emph{checkpoint}; nunca debe usarse \codigo{kill -9}, que cortaría la
escritura del \emph{checkpoint} a media subida.

\begin{lstlisting}[language=none,caption={Apagado graceful que preserva el checkpoint (node.service, recortado)}]
KillSignal=SIGTERM
TimeoutStopSec=120
\end{lstlisting}

El ciclo de vida del despliegue se cubre con scripts cortos para no gastar
cómputo cuando no se usa: la infraestructura normalmente está apagada y el día de
la revisión basta \codigo{deploy/encender.sh} para arrancar las cuatro VMs (los
tres nodos más el \codigo{generador}) y esperar a que el líder cargue las 820k
cuentas y sirva \codigo{/api/stats}, sondeándolo con reintentos de \codigo{curl}.
Al terminar, \codigo{deploy/apagar.sh} detiene las VMs sin borrarlas (pausa el
gasto) y \codigo{deploy/destruir-infra.sh} elimina todo, incluidos discos e IP
estática, para no dejar ningún costo. La IP estática del líder no cambia al
apagar y encender, así que el punto de entrada único es estable.

\subsection{Operacion el dia del concurso}

El guion de la demo en vivo está en \codigo{deploy/RUNBOOK.md}, organizado como
\textbf{encender $\rightarrow$ demostrar $\rightarrow$ apagar}. El paso 0 levanta
el clúster con \codigo{encender.sh} y abre el \emph{dashboard} en
\codigo{http://34.67.240.245:8080/} (la IP estática del líder). El paso 1 ejercita
a mano los cuatro \emph{endpoints} con JWT (\codigo{/api/register},
\codigo{/api/login}, \codigo{/api/accounts/\{id\}} y
\codigo{/api/transactions/transfer}) para mostrar sus formas exactas, y el paso 2
demuestra la invariante de conservación: \codigo{totalBalance} en
\codigo{/api/stats} debe quedar constante en \codigo{410113948069.86} antes y
después de transferir.

El corazón de la evaluación es la tolerancia a fallos (pasos 3 y 4). Se demuestra
el ``mata-proceso'' apagando primero \codigo{nodo-3} y luego \codigo{nodo-2} con
\codigo{gcloud compute instances stop}: el líder \textbf{sigue sirviendo} aunque
caigan las dos réplicas. Después se revive una réplica y se observa en su
\emph{log} que reanuda desde su secuencia \emph{exacta} (no desde cero), gracias
al \emph{checkpoint} que escribió al apagarse con SIGTERM, y que converge al líder
(mismo \codigo{lastTxId} y mismo \codigo{totalBalance}). El \emph{dashboard}
muestra todo el tiempo, por nodo, su estado, \#cuentas, saldo total,
\#transferencias, última tx y uso de CPU/RAM/disco, conservando los saldos.

\begin{lstlisting}[language=none,caption={Revivir una replica y verificar reanudacion por secuencia (RUNBOOK paso 4, recortado)}]
$GC compute instances start nodo-2 --zone=$ZONE     # revive
# en su log se ve que reanuda desde su secuencia EXACTA (no desde 0):
#   -> [checkpoint] restored checkpoint/nodo-2.json at watermark <N>
#   -> [node] replica resuming at sequence <N>
\end{lstlisting}

La medición del concurso se automatiza con \codigo{deploy/concurso.sh}, que corre
el \codigo{LoadDriver} \textbf{desde la VM generador} (baja latencia contra la IP
interna del líder) en tres escenarios encadenados, apagando una réplica entre uno
y otro: escenario 1 con el clúster completo (3 nodos), escenario 2 con líder $+$ 1
réplica y escenario 3 solo con el líder. Cada escenario verifica la invariante
(\codigo{CONSISTENT}) e imprime lecturas, transferencias y su
\codigo{score = transferencias*4 + lecturas}; el script suma los tres
\emph{scores} en un total. Que las tres corridas terminen \codigo{CONSISTENT} aun
cruzando apagados de réplicas es la prueba operativa de que el clúster no pierde
ni inventa dinero bajo fallos.

\begin{lstlisting}[language=none,caption={Un escenario del concurso sobre el generador (concurso.sh, recortado)}]
out="$(SSH_GEN "java -cp /opt/tesoreria-distribuida.jar mx.ipn.escom.tesoreria.loadtest.LoadDriver ${LEADER_INT} 8080 ${SECONDS_RUN} ${THREADS} 1 820000 ${label} 2>&1")"
echo "${out}" | grep -E 'CONSISTENT|INCONSISTENT|successful|score' >&2
\end{lstlisting}

Para cerrar, el RUNBOOK documenta la recuperación ante desastre total: como el
estado durable vive en Cloud Storage, un arranque en frío de los nodos (incluso
tras destruir las VMs y volver a crearlas) reconstruye el banco desde el bucket.
La limpieza final del día es \codigo{deploy/apagar.sh} para solo pausar el gasto,
o \codigo{deploy/destruir-infra.sh} para borrar discos e IP estática y no dejar
ningún costo. El RUNBOOK advierte además que las réplicas siempre se apaguen con
SIGTERM y nunca con \codigo{kill -9}, para no truncar la escritura del
\emph{checkpoint}.


\end{document}
